\input{preamble}

% OK, commencer ici.
%
\begin{document}

\title{Problèmes de déformation}

\maketitle

\phantomsection
\label{section-phantom}

\tableofcontents

\section{Introduction}
\label{section-introduction}

\noindent
Le but de ce chapitre est de traiter des exemples de la théorie générale
développée dans les chapitres Théorie formelle des déformations,
Théorie des déformations et Le complexe cotangent.

\medskip\noindent
La section 3 de l'article \cite{Sch} de Schlessinger étudie elle aussi
quelques exemples.






\section{Exemples de problèmes de déformation}
\label{section-examples}

\noindent
Liste des sujets qui devraient figurer ici :
\begin{enumerate}
\item Déformations de schémas :
\begin{enumerate}
\item La condition de Rim-Schlessinger.
\item Calcul de l'espace tangent.
\item Calcul des déformations infinitésimales.
\item La catégorie de déformations d'une hypersurface affine.
\end{enumerate}
\item Déformations de faisceaux (par exemple, fixer $X/S$, un point $s$
de type fini de $S$ et un faisceau quasi-cohérent $\mathcal{F}_s$ sur $X_s$).
\item Déformations d'espaces algébriques (très semblables aux déformations
de schémas ; peut-être même plus faciles ?).
\item Déformations d'applications (par exemple de morphismes de schémas ; on peut fixer
à la fois la source et le but, ou seulement l'un des deux).
\item Ajouter d'autres sujets ici.
\end{enumerate}





\section{Plan général}
\label{section-general}

\noindent
Cette section expose la démarche suivie pour étudier les exemples qui viennent.

\medskip\noindent
Étape I. Dans chaque section, on fixe un anneau noethérien $\Lambda$ et
un morphisme fini d'anneaux $\Lambda \to k$, où $k$ est un corps.
Comme d'habitude, on prend $\mathcal{C}_\Lambda = \mathcal{C}_{\Lambda, k}$
pour catégorie de base ; voir Théorie formelle des déformations,
définition \ref{formal-defos-definition-CLambda}.

\medskip\noindent
Étape II. Dans chaque section, on définit une catégorie $\mathcal{F}$
cofibrée en groupoïdes sur $\mathcal{C}_\Lambda$. Il nous arrivera
de considérer à la place un foncteur
$F : \mathcal{C}_\Lambda \to \textit{Ensembles}$.

\medskip\noindent
Étape III. On explique dans quelle mesure $\mathcal{F}$ satisfait
la condition de Rim-Schlessinger (RS) étudiée dans Théorie formelle
des déformations, section \ref{formal-defos-section-RS-condition}.
De même, on pourra étudier dans quelle mesure $\mathcal{F}$
satisfait (S1) et (S2), ou dans quelle mesure $F$ satisfait
les conditions correspondantes (H1) et (H2) de Schlessinger.
Voir Théorie formelle des déformations, section
\ref{formal-defos-section-schlessinger-conditions}.

\medskip\noindent
Étape IV. Soit $x_0$ un objet de $\mathcal{F}(k)$, autrement dit un objet
de $\mathcal{F}$ sur $k$. Dans ce chapitre, on emploiera la notation
$$
\Deformationcategory_{x_0} = \mathcal{F}_{x_0}
$$
pour désigner la catégorie de prédéformations construite dans Théorie
formelle des déformations, remarque
\ref{formal-defos-remark-localize-cofibered-groupoid}.
Si $\mathcal{F}$ satisfait (RS), alors
$\Deformationcategory_{x_0}$ est une catégorie de déformations
(Théorie formelle des déformations, lemme
\ref{formal-defos-lemma-localize-RS})
et satisfait (S1) et (S2)
(Théorie formelle des déformations, lemme
\ref{formal-defos-lemma-RS-implies-S1-S2}).
Si (S1) et (S2) sont satisfaites, une question importante
est de savoir si l'espace tangent
$$
T\Deformationcategory_{x_0} = T_{x_0}\mathcal{F} = T\mathcal{F}_{x_0}
$$
(voir Théorie formelle des déformations, remarque
\ref{formal-defos-remark-tangent-space-cofibered-groupoid} et
définition \ref{formal-defos-definition-tangent-space})
est de dimension finie. En effet, cela garantit que
$\Deformationcategory_{x_0}$ possède un objet formel versel
(Théorie formelle des déformations, lemme
\ref{formal-defos-lemma-versal-object-existence}).

\medskip\noindent
Étape V. Si $\mathcal{F}$ franchit l'étape IV, la question suivante est de savoir
si le $k$-espace vectoriel
$$
\text{Inf}(\Deformationcategory_{x_0}) = \text{Inf}_{x_0}(\mathcal{F})
$$
des automorphismes infinitésimaux de $x_0$ est de dimension finie.
Dans l'affirmative, cela implique que
$\Deformationcategory_{x_0}$ admet une présentation par un
groupoïde lisse proreprésentable en foncteurs sur $\mathcal{C}_\Lambda$ ; voir
Théorie formelle des déformations, théorème
\ref{formal-defos-theorem-presentation-deformation-groupoid}.





\section{Modules projectifs de type fini}
\label{section-finite-projective-modules}

\noindent
Cette section n'est qu'un échauffement. Bien entendu, les modules projectifs
de type fini ne devraient avoir aucun « module » de paramètres.

\begin{example}[Modules projectifs de type fini]
\label{example-finite-projective-modules}
Soit $\mathcal{F}$ la catégorie définie comme suit :
\begin{enumerate}
\item un objet est un couple $(A, M)$ formé d'un
objet $A$ de $\mathcal{C}_\Lambda$ et d'un $A$-module
projectif de type fini $M$ ;
\item un morphisme $(f, g) : (B, N) \to (A, M)$ est formé d'un
morphisme $f : B \to A$ dans $\mathcal{C}_\Lambda$ et
d'une application $g : N \to M$ qui est $f$-linéaire et induit
un isomorphisme $N \otimes_{B, f} A \cong M$.
\end{enumerate}
Le foncteur $p : \mathcal{F} \to \mathcal{C}_\Lambda$ envoie $(A, M)$ sur $A$
et $(f, g)$ sur $f$. Il est clair que $p$ est cofibré en groupoïdes.
Étant donné un $k$-espace vectoriel de dimension finie $V$,
soit $x_0 = (k, V)$ l'objet correspondant de $\mathcal{F}(k)$.
On pose
$$
\Deformationcategory_V = \mathcal{F}_{x_0}
$$
\end{example}

\noindent
Puisque tout module projectif de type fini sur un anneau local est libre
de type fini (Algèbre, lemme \ref{algebra-lemma-finite-projective}), on a
$$
\begin{matrix}
\text{classes d'isomorphisme} \\
\text{des objets de }\mathcal{F}(A)
\end{matrix}
= \coprod\nolimits_{n \geq 0} \{*\}
$$
Bien que cela signifie que la théorie des déformations de $\mathcal{F}$
est essentiellement triviale, on parcourt néanmoins les étapes décrites
à la section \ref{section-general} afin de donner un exemple simple.

\begin{lemma}
\label{lemma-finite-projective-modules-RS}
L'exemple \ref{example-finite-projective-modules}
satisfait la condition de Rim-Schlessinger (RS).
En particulier, $\Deformationcategory_V$ est une catégorie de déformations
pour tout espace vectoriel $V$ de dimension finie sur $k$.
\end{lemma}

\begin{proof}
Soient $A_1 \to A$ et $A_2 \to A$ des morphismes de $\mathcal{C}_\Lambda$.
Supposons $A_2 \to A$ surjectif. D'après Théorie formelle
des déformations, lemme
\ref{formal-defos-lemma-RS-2-categorical}
il suffit de montrer que le foncteur
$\mathcal{F}(A_1 \times_A A_2) \to
\mathcal{F}(A_1) \times_{\mathcal{F}(A)} \mathcal{F}(A_2)$
est une équivalence de catégories.

\medskip\noindent
Il faut donc montrer que la catégorie des modules projectifs de type fini
sur $A_1 \times_A A_2$ est équivalente au produit fibré
des catégories des modules projectifs de type fini sur $A_1$ et $A_2$
au-dessus de la catégorie des modules projectifs de type fini sur $A$.
C'est un cas particulier de Compléments d'algèbre, lemme
\ref{more-algebra-lemma-finitely-presented-module-over-fibre-product}.
Rappelons que le foncteur inverse envoie le triplet
$(M_1, M_2, \varphi)$, où
$M_1$ est un $A_1$-module projectif de type fini,
$M_2$ est un $A_2$-module projectif de type fini et
$\varphi : M_1 \otimes_{A_1} A \to M_2 \otimes_{A_2} A$
est un isomorphisme de $A$-modules, sur le $A_1 \times_A A_2$-module
projectif de type fini $M_1 \times_\varphi M_2$.
\end{proof}

\begin{lemma}
\label{lemma-finite-projective-modules-TI}
Dans l'exemple \ref{example-finite-projective-modules},
soit $V$ un $k$-espace vectoriel de dimension finie. Alors
$$
T\Deformationcategory_V = (0)
\quad\text{et}\quad
\text{Inf}(\Deformationcategory_V) = \text{End}_k(V)
$$
sont de dimension finie.
\end{lemma}

\begin{proof}
Avec $\mathcal{F}$ comme dans l'exemple \ref{example-finite-projective-modules},
posons $x_0 = (k, V) \in \Ob(\mathcal{F}(k))$.
Rappelons que $T\Deformationcategory_V = T_{x_0}\mathcal{F}$
est l'ensemble des classes d'isomorphisme
de couples $(x, \alpha)$ formés d'un objet $x$ de $\mathcal{F}$
sur l'anneau des nombres duaux $k[\epsilon]$ et d'un morphisme
$\alpha : x \to x_0$ de $\mathcal{F}$ au-dessus de $k[\epsilon] \to k$.

\medskip\noindent
À isomorphisme près, il existe un unique couple $(M, \alpha)$ formé d'un
module projectif de type fini $M$ sur $k[\epsilon]$
et d'une application $k[\epsilon]$-linéaire $\alpha : M \to V$
qui induit un isomorphisme $M \otimes_{k[\epsilon]} k \to V$.
Par exemple, si $V = k^{\oplus n}$, on prend
$M = k[\epsilon]^{\oplus n}$ muni de l'application évidente $\alpha$.

\medskip\noindent
De même, $\text{Inf}(\Deformationcategory_V) = \text{Inf}_{x_0}(\mathcal{F})$
est l'ensemble des automorphismes
de la déformation triviale $x'_0$ de $x_0$ sur $k[\epsilon]$.
Voir Théorie formelle des déformations, définition
\ref{formal-defos-definition-infinitesimal-auts} pour plus de détails.

\medskip\noindent
Pour $(M, \alpha)$ comme au deuxième paragraphe, on voit qu'un élément de
$\text{Inf}_{x_0}(\mathcal{F})$ est un automorphisme $\gamma : M \to M$ tel que
$\gamma \bmod \epsilon = \text{id}$. On peut donc écrire
$\gamma = \text{id}_M + \epsilon \psi$, où
$\psi : M/\epsilon M \to M/\epsilon M$ est $k$-linéaire.
Au moyen de $\alpha$, on peut voir $\psi$ comme un élément de
$\text{End}_k(V)$, ce qui achève la démonstration.
\end{proof}



\section{Représentations d'un groupe}
\label{section-representations}

\noindent
La théorie des déformations des représentations peut être très intéressante.

\begin{example}[Représentations d'un groupe]
\label{example-representations}
Soit $\Gamma$ un groupe.
Soit $\mathcal{F}$ la catégorie définie comme suit :
\begin{enumerate}
\item un objet est un triplet $(A, M, \rho)$ formé d'un
objet $A$ de $\mathcal{C}_\Lambda$, d'un $A$-module projectif de type fini $M$
et d'un homomorphisme $\rho : \Gamma \to \text{GL}_A(M)$ ;
\item un morphisme $(f, g) : (B, N, \tau) \to (A, M, \rho)$ est formé d'un
morphisme $f : B \to A$ dans $\mathcal{C}_\Lambda$ et
d'une application $g : N \to M$ qui est $f$-linéaire, $\Gamma$-équivariante
et induit un isomorphisme $N \otimes_{B, f} A \cong M$.
\end{enumerate}
Le foncteur $p : \mathcal{F} \to \mathcal{C}_\Lambda$ envoie $(A, M, \rho)$
sur $A$ et $(f, g)$ sur $f$. Il est clair que $p$ est cofibré en groupoïdes.
Étant donné un $k$-espace vectoriel de dimension finie $V$ et une représentation
$\rho_0 : \Gamma \to \text{GL}_k(V)$,
soit $x_0 = (k, V, \rho_0)$ l'objet correspondant de $\mathcal{F}(k)$.
On pose
$$
\Deformationcategory_{V, \rho_0} = \mathcal{F}_{x_0}
$$
\end{example}

\noindent
Puisque tout module projectif de type fini sur un anneau local est libre
de type fini (Algèbre, lemme \ref{algebra-lemma-finite-projective}),
on a
$$
\begin{matrix}
\text{classes d'isomorphisme} \\
\text{des objets de }\mathcal{F}(A)
\end{matrix}
=
\coprod\nolimits_{n \geq 0}\quad
\begin{matrix}
\text{GL}_n(A)\text{-orbites de conjugaison des}\\
\text{homomorphismes }\rho : \Gamma \to \text{GL}_n(A)
\end{matrix}
$$
C'est déjà plus intéressant que l'étude menée à la
section \ref{section-finite-projective-modules}.

\begin{lemma}
\label{lemma-representations-RS}
L'exemple \ref{example-representations}
satisfait la condition de Rim-Schlessinger (RS).
En particulier, $\Deformationcategory_{V, \rho_0}$ est une catégorie de déformations
pour toute représentation de dimension finie
$\rho_0 : \Gamma \to \text{GL}_k(V)$.
\end{lemma}

\begin{proof}
Soient $A_1 \to A$ et $A_2 \to A$ des morphismes de $\mathcal{C}_\Lambda$.
Supposons $A_2 \to A$ surjectif. D'après Théorie formelle
des déformations, lemme
\ref{formal-defos-lemma-RS-2-categorical}
il suffit de montrer que le foncteur
$\mathcal{F}(A_1 \times_A A_2) \to
\mathcal{F}(A_1) \times_{\mathcal{F}(A)} \mathcal{F}(A_2)$
est une équivalence de catégories.

\medskip\noindent
Considérons un objet
$$
((A_1, M_1, \rho_1), (A_2, M_2, \rho_2), (\text{id}_A, \varphi))
$$
de la catégorie $\mathcal{F}(A_1) \times_{\mathcal{F}(A)} \mathcal{F}(A_2)$.
Alors, comme on l'a vu dans la démonstration du lemme \ref{lemma-finite-projective-modules-RS},
on peut considérer le $A_1 \times_A A_2$-module projectif
de type fini $M_1 \times_\varphi M_2$.
Puisque $\varphi$ est compatible aux actions données, on obtient
$$
\rho_1 \times \rho_2 : \Gamma \longrightarrow
\text{GL}_{A_1 \times_A A_2}(M_1 \times_\varphi M_2)
$$
Alors $(M_1 \times_\varphi M_2, \rho_1 \times \rho_2)$
est un objet de $\mathcal{F}(A_1 \times_A A_2)$.
Cette construction détermine un quasi-inverse de notre foncteur.
\end{proof}

\begin{lemma}
\label{lemma-representations-TI}
Dans l'exemple \ref{example-representations}, soit
$\rho_0 : \Gamma \to \text{GL}_k(V)$
une représentation de dimension finie. Alors
$$
T\Deformationcategory_{V, \rho_0} = \Ext^1_{k[\Gamma]}(V, V) =
H^1(\Gamma, \text{End}_k(V))
\quad\text{et}\quad
\text{Inf}(\Deformationcategory_{V, \rho_0}) = H^0(\Gamma, \text{End}_k(V))
$$
Ainsi, $\text{Inf}(\Deformationcategory_{V, \rho_0})$
est toujours de dimension finie,
et $T\Deformationcategory_{V, \rho_0}$ est de dimension finie
si $\Gamma$ est de type fini.
\end{lemma}

\begin{proof}
Traitons d'abord les automorphismes infinitésimaux.
Soit $M = V \otimes_k k[\epsilon]$, muni de l'action induite
$\rho_0' : \Gamma \to \text{GL}_n(M)$.
Un automorphisme infinitésimal, c'est-à-dire un élément de
$\text{Inf}(\Deformationcategory_{V, \rho_0})$,
est alors donné par un automorphisme
$\gamma = \text{id} + \epsilon \psi : M \to M$
comme dans la démonstration du lemme \ref{lemma-finite-projective-modules-TI},
où, de plus, $\psi$ doit commuter
à l'action de $\Gamma$ (donnée par $\rho_0$).
On voit donc que
$$
\text{Inf}(\Deformationcategory_{V, \rho_0}) = H^0(\Gamma, \text{End}_k(V))
$$
comme l'affirme le lemme.

\medskip\noindent
Ensuite, soit $(k[\epsilon], M, \rho)$ un objet de $\mathcal{F}$
sur $k[\epsilon]$, et soit $\alpha : M \to V$ une application $\Gamma$-équivariante
induisant un isomorphisme $M/\epsilon M \to V$.
Puisque $M$ est libre comme $k[\epsilon]$-module, on obtient
une extension de $\Gamma$-modules
$$
0 \to V \to M \xrightarrow{\alpha} V \to 0
$$
On omet la construction détaillée de l'application de gauche.
Réciproquement, si l'on dispose d'une extension de $\Gamma$-modules comme
ci-dessus, on peut l'utiliser pour définir une structure de $k[\epsilon]$-module
sur $M$ et obtenir un objet de $\mathcal{F}(k[\epsilon])$
ainsi qu'une application $\alpha$ comme ci-dessus.
Il s'ensuit que
$$
T\Deformationcategory_{V, \rho_0} = \Ext^1_{k[\Gamma]}(V, V)
$$
comme l'affirme le lemme. Cet espace est égal à
$H^1(\Gamma, \text{End}_k(V))$ d'après Cohomologie étale,
lemme \ref{etale-cohomology-lemma-ext-modules-hom}.

\medskip\noindent
L'assertion sur les dimensions résulte de Cohomologie étale,
lemme
\ref{etale-cohomology-lemma-finite-dim-group-cohomology}.
\end{proof}

\noindent
Dans l'exemple \ref{example-representations}, si $\Gamma$ est de type fini
et si $(V, \rho_0)$ est une représentation de dimension finie de $\Gamma$
sur $k$, alors $\Deformationcategory_{V, \rho_0}$
admet une présentation par un groupoïde lisse proreprésentable en foncteurs
sur $\mathcal{C}_\Lambda$
et possède a fortiori un objet formel versel (minimal). Cela résulte
des lemmes \ref{lemma-representations-RS} et \ref{lemma-representations-TI}
ainsi que de la discussion générale de la section \ref{section-general}.

\begin{lemma}
\label{lemma-representations-hull}
Dans l'exemple \ref{example-representations}, supposons $\Gamma$ de type fini.
Soit $\rho_0 : \Gamma \to \text{GL}_k(V)$ une représentation de dimension finie.
Supposons que $\Lambda$ soit un anneau local complet de corps résiduel $k$
(cas classique). Alors le foncteur
$$
F : \mathcal{C}_\Lambda \longrightarrow \textit{Ensembles},\quad
A \longmapsto \Ob(\Deformationcategory_{V, \rho_0}(A))/\cong
$$
des classes d'isomorphisme d'objets possède une enveloppe. Si
$H^0(\Gamma, \text{End}_k(V)) = k$, alors $F$ est
proreprésentable.
\end{lemma}

\begin{proof}
L'existence d'une enveloppe résulte des lemmes \ref{lemma-representations-RS} et
\ref{lemma-representations-TI} et
Théorie formelle des déformations, lemme \ref{formal-defos-lemma-RS-implies-S1-S2}
et remarque \ref{formal-defos-remark-compose-minimal-into-iso-classes}.

\medskip\noindent
Supposons $H^0(\Gamma, \text{End}_k(V)) = k$. Pour voir que $F$
est proreprésentable, il suffit de montrer que $F$ est un
foncteur de déformations ; voir Théorie formelle des déformations, théorème
\ref{formal-defos-theorem-Schlessinger-prorepresentability}.
Autrement dit, il faut montrer que $F$ satisfait (RS).
On peut pour cela utiliser le critère de Théorie formelle des déformations, lemme
\ref{formal-defos-lemma-RS-associated-functor}.
La surjectivité requise des groupes d'automorphismes résultera de ce que
l'on montre que
$$
A \cdot \text{id}_M =
\text{End}_{A[\Gamma]}(M)
$$
pour tout objet $(A, M, \rho)$ de $\mathcal{F}$ tel que
$M \otimes_A k$ soit isomorphe à $V$ comme représentation de $\Gamma$.
Puisque le membre de gauche est contenu dans celui de droite,
il suffit de montrer
$\text{length}_A \text{End}_{A[\Gamma]}(M) \leq \text{length}_A A$.
Choisissons des idéaux deux à deux distincts
$(0) = I_n \subset \ldots \subset I_1 \subset A$
avec $n = \text{length}(A)$. En filtrant $M$
de façon correspondante, on voit qu'il suffit de prouver que $\Hom_{A[\Gamma]}(M, I_tM/I_{t + 1}M)$
est de longueur $1$. Puisque $I_tM/I_{t + 1}M \cong M \otimes_A k$
et que toute application de $A[\Gamma]$-modules $M \to M \otimes_A k$ se factorise
de manière unique par l'application quotient $M \to M \otimes_A k$
et donne ainsi un élément de
$$
\text{End}_{A[\Gamma]}(M \otimes_A k) = \text{End}_{k[\Gamma]}(V) = k
$$
on conclut.
\end{proof}



\section{Représentations continues}
\label{section-continuous-representations}

\noindent
Une démarche particulièrement intéressante consiste à prendre un groupe de Galois
infini et à étudier la théorie des déformations de ses représentations ; voir
\cite{Mazur-deforming}.

\begin{example}[Représentations d'un groupe topologique]
\label{example-continuous-representations}
Soit $\Gamma$ un groupe topologique.
Soit $\mathcal{F}$ la catégorie définie comme suit :
\begin{enumerate}
\item un objet est un triplet $(A, M, \rho)$ formé d'un
objet $A$ de $\mathcal{C}_\Lambda$, d'un $A$-module projectif de type fini $M$
et d'un homomorphisme continu $\rho : \Gamma \to \text{GL}_A(M)$,
où $\text{GL}_A(M)$ est muni de la topologie discrète\footnote{Une autre possibilité
serait d'exiger que le $A$-module $M$, muni de l'action de $G$ donnée par $\rho$,
soit un $A\text{-}G$-module au sens de Cohomologie étale, définition
\ref{etale-cohomology-definition-G-module-continuous}. Cependant,
puisque $M$ est un $A$-module de type fini, ces conditions sont équivalentes.} ;
\item un morphisme $(f, g) : (B, N, \tau) \to (A, M, \rho)$ est formé d'un
morphisme $f : B \to A$ dans $\mathcal{C}_\Lambda$ et
d'une application $g : N \to M$ qui est $f$-linéaire, $\Gamma$-équivariante
et induit un isomorphisme $N \otimes_{B, f} A \cong M$.
\end{enumerate}
Le foncteur $p : \mathcal{F} \to \mathcal{C}_\Lambda$ envoie $(A, M, \rho)$
sur $A$ et $(f, g)$ sur $f$. Il est clair que $p$ est cofibré en groupoïdes.
Étant donné un $k$-espace vectoriel de dimension finie $V$ et une
représentation continue $\rho_0 : \Gamma \to \text{GL}_k(V)$,
soit $x_0 = (k, V, \rho_0)$ l'objet correspondant de $\mathcal{F}(k)$.
On pose
$$
\Deformationcategory_{V, \rho_0} = \mathcal{F}_{x_0}
$$
\end{example}

\noindent
Puisque tout module projectif de type fini sur un anneau local est libre
de type fini (Algèbre, lemme \ref{algebra-lemma-finite-projective}),
on a
$$
\begin{matrix}
\text{classes d'isomorphisme} \\
\text{des objets de }\mathcal{F}(A)
\end{matrix}
=
\coprod\nolimits_{n \geq 0}\quad
\begin{matrix}
\text{GL}_n(A)\text{-orbites de conjugaison des}\\
\text{homomorphismes continus }\rho : \Gamma \to \text{GL}_n(A)
\end{matrix}
$$

\begin{lemma}
\label{lemma-continuous-representations-RS}
L'exemple \ref{example-continuous-representations}
satisfait la condition de Rim-Schlessinger (RS).
En particulier, $\Deformationcategory_{V, \rho_0}$ est une catégorie de déformations
pour toute représentation continue de dimension finie
$\rho_0 : \Gamma \to \text{GL}_k(V)$.
\end{lemma}

\begin{proof}
La démonstration est exactement la même que celle du
lemme \ref{lemma-representations-RS}.
\end{proof}

\begin{lemma}
\label{lemma-continuous-representations-TI}
Dans l'exemple \ref{example-continuous-representations}, soit
$\rho_0 : \Gamma \to \text{GL}_k(V)$ une représentation continue
de dimension finie. Alors
$$
T\Deformationcategory_{V, \rho_0} = H^1(\Gamma, \text{End}_k(V))
\quad\text{et}\quad
\text{Inf}(\Deformationcategory_{V, \rho_0}) = H^0(\Gamma, \text{End}_k(V))
$$
Ainsi, $\text{Inf}(\Deformationcategory_{V, \rho_0})$
est toujours de dimension finie,
et $T\Deformationcategory_{V, \rho_0}$ est de dimension finie
si $\Gamma$ est topologiquement de type fini.
\end{lemma}

\begin{proof}
La démonstration est exactement la même que celle du
lemme \ref{lemma-representations-TI}.
\end{proof}

\noindent
Dans l'exemple \ref{example-continuous-representations}, si $\Gamma$
est topologiquement de type fini
et si $(V, \rho_0)$ est une représentation continue de dimension finie de $\Gamma$
sur $k$, alors $\Deformationcategory_{V, \rho_0}$
admet une présentation par un groupoïde lisse proreprésentable en foncteurs
sur $\mathcal{C}_\Lambda$
et possède a fortiori un objet formel versel (minimal). Cela résulte
des lemmes \ref{lemma-continuous-representations-RS} et
\ref{lemma-continuous-representations-TI}
ainsi que de la discussion générale de la section \ref{section-general}.

\begin{lemma}
\label{lemma-continuous-representations-hull}
Dans l'exemple \ref{example-continuous-representations}, supposons que $\Gamma$
soit topologiquement de type fini.
Soit $\rho_0 : \Gamma \to \text{GL}_k(V)$ une représentation de dimension finie.
Supposons que $\Lambda$ soit un anneau local complet de corps résiduel $k$
(cas classique). Alors le foncteur
$$
F : \mathcal{C}_\Lambda \longrightarrow \textit{Ensembles},\quad
A \longmapsto \Ob(\Deformationcategory_{V, \rho_0}(A))/\cong
$$
des classes d'isomorphisme d'objets possède une enveloppe. Si
$H^0(\Gamma, \text{End}_k(V)) = k$, alors $F$ est
proreprésentable.
\end{lemma}

\begin{proof}
La démonstration est exactement la même que celle du
lemme \ref{lemma-representations-hull}.
\end{proof}



\section{Algèbres graduées}
\label{section-graded-algebras}

\noindent
Nous utiliserons l'exemple de cette section dans la démonstration que le champ des
schémas propres polarisés est un champ algébrique. Pour cette raison, nous
considérerons des algèbres graduées commutatives dont les composantes homogènes sont
des modules projectifs de type fini (parfois dits « localement finis »).

\begin{example}[Algèbres graduées]
\label{example-graded-algebras}
Soit $\mathcal{F}$ la catégorie définie comme suit :
\begin{enumerate}
\item un objet est un couple $(A, P)$ formé d'un
objet $A$ de $\mathcal{C}_\Lambda$ et d'une $A$-algèbre graduée $P$
tel que $P_d$ soit un $A$-module projectif de type fini pour tout $d \geq 0$ ;
\item un morphisme $(f, g) : (B, Q) \to (A, P)$ est formé d'un
morphisme $f : B \to A$ dans $\mathcal{C}_\Lambda$ et
d'une application $g : Q \to P$ qui est $f$-linéaire et induit un
isomorphisme $Q \otimes_{B, f} A \cong P$.
\end{enumerate}
Le foncteur $p : \mathcal{F} \to \mathcal{C}_\Lambda$ envoie $(A, P)$
sur $A$ et $(f, g)$ sur $f$. Il est clair que $p$ est cofibré en groupoïdes.
Étant donnée une $k$-algèbre graduée $P$ telle que $\dim_k(P_d) < \infty$ pour tout
$d \geq 0$, soit $x_0 = (k, P)$ l'objet correspondant de $\mathcal{F}(k)$.
On pose
$$
\Deformationcategory_P = \mathcal{F}_{x_0}
$$
\end{example}

\begin{lemma}
\label{lemma-graded-algebras-RS}
L'exemple \ref{example-graded-algebras}
satisfait la condition de Rim-Schlessinger (RS).
En particulier, $\Deformationcategory_P$ est une catégorie de déformations
pour toute $k$-algèbre graduée $P$.
\end{lemma}

\begin{proof}
Soient $A_1 \to A$ et $A_2 \to A$ des morphismes de $\mathcal{C}_\Lambda$.
Supposons $A_2 \to A$ surjectif. D'après Théorie formelle
des déformations, lemme
\ref{formal-defos-lemma-RS-2-categorical}
il suffit de montrer que le foncteur
$\mathcal{F}(A_1 \times_A A_2) \to
\mathcal{F}(A_1) \times_{\mathcal{F}(A)} \mathcal{F}(A_2)$
est une équivalence de catégories.

\medskip\noindent
Considérons un objet
$$
((A_1, P_1), (A_2, P_2), (\text{id}_A, \varphi))
$$
de la catégorie $\mathcal{F}(A_1) \times_{\mathcal{F}(A)} \mathcal{F}(A_2)$.
Considérons alors $P_1 \times_\varphi P_2$. Comme
$\varphi : P_1 \otimes_{A_1} A \to P_2 \otimes_{A_2} A$
est un isomorphisme d'algèbres graduées, on voit que les composantes homogènes
de $P_1 \times_\varphi P_2$ sont des $A_1 \times_A A_2$-modules projectifs de type fini ;
voir la démonstration du lemme \ref{lemma-finite-projective-modules-RS}.
Ainsi, $P_1 \times_\varphi P_2$ est un objet de $\mathcal{F}(A_1 \times_A A_2)$.
Cette construction détermine un quasi-inverse de notre foncteur,
ce qui achève la démonstration.
\end{proof}

\begin{lemma}
\label{lemma-graded-algebras-TI}
Dans l'exemple \ref{example-graded-algebras}, soit $P$ une $k$-algèbre graduée.
Alors
$$
T\Deformationcategory_P
\quad\text{et}\quad
\text{Inf}(\Deformationcategory_P) = \text{Der}_k(P, P)
$$
sont de dimension finie si $P$ est de type fini sur $k$.
\end{lemma}

\begin{proof}
Traitons d'abord les automorphismes infinitésimaux.
Soit $Q = P \otimes_k k[\epsilon]$.
Un élément de $\text{Inf}(\Deformationcategory_P)$
est alors donné par un automorphisme
$\gamma = \text{id} + \epsilon \delta : Q \to Q$
comme ci-dessus, où maintenant $\delta : P \to P$.
Le fait que $\gamma$ soit gradué implique que
$\delta$ est homogène de degré $0$.
Le fait que $\gamma$ soit $k$-linéaire implique que
$\delta$ est $k$-linéaire.
Le fait que $\gamma$ soit multiplicatif implique que
$\delta$ est une $k$-dérivation.
Réciproquement, toute $k$-dérivation $\delta : P \to P$
homogène de degré $0$ fournit un automorphisme
$\gamma = \text{id} + \epsilon \delta$ comme ci-dessus.
On voit donc que
$$
\text{Inf}(\Deformationcategory_P) = \text{Der}_k(P, P)
$$
comme l'affirme le lemme.
Manifestement, si $P$ est engendrée par les composantes $P_i$ de degrés
$0 \leq i \leq N$, alors $\delta$ est déterminée par
les applications linéaires $\delta_i : P_i \to P_i$ pour
$0 \leq i \leq N$, et l'on voit que
$$
\dim_k \text{Der}_k(P, P) < \infty
$$
comme voulu.

\medskip\noindent
Pour achever la démonstration du lemme, montrons que l'espace des déformations
est de dimension finie. À cette fin, choisissons
une présentation
$$
k[X_1, \ldots, X_n]/(F_1, \ldots, F_m) \longrightarrow P
$$
de $k$-algèbres graduées, où $\deg(X_i) = d_i$ et où
$F_j$ est homogène de degré $e_j$.
Soit $Q$ une $k[\epsilon]$-algèbre graduée quelconque,
libre de type fini en chaque degré, munie d'un isomorphisme
$\alpha : Q/\epsilon Q \to P$ tel que $(Q, \alpha)$ définisse
un élément de $T\Deformationcategory_P$.
Choisissons un élément homogène $q_i \in Q$ de degré $d_i$
dont l'image est celle de $X_i$ dans $P$.
On obtient alors
$$
k[\epsilon][X_1, \ldots, X_n] \longrightarrow Q,\quad
X_i \longmapsto q_i
$$
et, puisque $P = Q/\epsilon Q$, cette application est surjective d'après le lemme de Nakayama.
Une chasse au diagramme élémentaire montre que l'on peut choisir des éléments homogènes
$F_{\epsilon, j} \in k[\epsilon][X_1, \ldots, X_n]$ de degré $e_j$
dont l'image dans $Q$ est nulle et dont l'image est $F_j$ dans $k[X_1, \ldots, X_n]$.
Alors
$$
k[\epsilon][X_1, \ldots, X_n]/(F_{\epsilon, 1}, \ldots, F_{\epsilon, m})
\longrightarrow Q
$$
est une présentation de $Q$, par platitude de $Q$ sur $k[\epsilon]$.
Écrivons
$$
F_{\epsilon, j} =  F_j + \epsilon G_j
$$
Le vecteur $(G_1, \ldots, G_m)$ comporte une certaine ambiguïté.
D'abord, en choisissant autrement les $F_{\epsilon, j}$,
on peut modifier $G_j$ par un élément arbitraire de degré $e_j$
du noyau de $k[X_1, \ldots, X_n] \to P$.
Ainsi, au lieu de $(G_1, \ldots, G_m)$, nous retenons
l'élément
$$
(g_1, \ldots, g_m) \in P_{e_1} \oplus \ldots \oplus P_{e_m}
$$
où $g_j$ est l'image de $G_j$ dans $P_{e_j}$.
En outre, si l'on remplace le choix de $q_i$ par $q_i + \epsilon p_i$,
où $p_i$ est de degré $d_i$, un calcul (omis) montre
que $g_j$ devient
$$
g_j^{new} = g_j - \sum\nolimits_{i = 1}^n p_i \partial F_j / \partial X_i
$$
On conclut que la classe d'isomorphisme de $Q$ est déterminée par
l'image du vecteur $(G_1, \ldots, G_m)$ dans le $k$-espace vectoriel
$$
W  = \Coker(P_{d_1} \oplus \ldots \oplus P_{d_n}
\xrightarrow{(\frac{\partial F_j}{\partial X_i})}
P_{e_1} \oplus \ldots \oplus P_{e_m})
$$
On obtient ainsi une injection
$$
T\Deformationcategory_P \longrightarrow W
$$
Puisque $W$ est manifestement de dimension finie, le lemme en résulte.
\end{proof}

\noindent
Dans l'exemple \ref{example-graded-algebras}, si $P$ est une $k$-algèbre graduée
de type fini, alors $\Deformationcategory_P$
admet une présentation par un groupoïde lisse proreprésentable en foncteurs
sur $\mathcal{C}_\Lambda$
et possède a fortiori un objet formel versel (minimal). Cela résulte
des lemmes \ref{lemma-graded-algebras-RS} et
\ref{lemma-graded-algebras-TI}
ainsi que de la discussion générale de la section \ref{section-general}.

\begin{lemma}
\label{lemma-graded-algebras-hull}
Dans l'exemple \ref{example-graded-algebras}, supposons que $P$ soit une
$k$-algèbre graduée de type fini. Supposons que $\Lambda$ soit un anneau local complet
de corps résiduel $k$
(cas classique). Alors le foncteur
$$
F : \mathcal{C}_\Lambda \longrightarrow \textit{Ensembles},\quad
A \longmapsto \Ob(\Deformationcategory_P(A))/\cong
$$
des classes d'isomorphisme d'objets possède une enveloppe.
\end{lemma}

\begin{proof}
Cela résulte immédiatement des lemmes \ref{lemma-graded-algebras-RS} et
\ref{lemma-graded-algebras-TI} ainsi que de
Théorie formelle des déformations, lemme \ref{formal-defos-lemma-RS-implies-S1-S2}
et remarque \ref{formal-defos-remark-compose-minimal-into-iso-classes}.
\end{proof}







\section{Anneaux}
\label{section-rings}

\noindent
La théorie des déformations des anneaux est la même que celle des
schémas affines. Pour les anneaux et les schémas, lorsque nous parlons de déformations,
nous entendons des déformations {\it plates}.

\begin{example}[Anneaux]
\label{example-rings}
Soit $\mathcal{F}$ la catégorie définie comme suit :
\begin{enumerate}
\item un objet est un couple $(A, P)$ formé d'un
objet $A$ de $\mathcal{C}_\Lambda$ et d'une $A$-algèbre plate $P$ ;
\item un morphisme $(f, g) : (B, Q) \to (A, P)$ est formé d'un
morphisme $f : B \to A$ dans $\mathcal{C}_\Lambda$ et
d'une application $g : Q \to P$ qui est $f$-linéaire et induit un
isomorphisme $Q \otimes_{B, f} A \cong P$.
\end{enumerate}
Le foncteur $p : \mathcal{F} \to \mathcal{C}_\Lambda$ envoie $(A, P)$
sur $A$ et $(f, g)$ sur $f$. Il est clair que $p$ est cofibré en groupoïdes.
Étant donnée une $k$-algèbre $P$, soit $x_0 = (k, P)$ l'objet correspondant
de $\mathcal{F}(k)$. On pose
$$
\Deformationcategory_P = \mathcal{F}_{x_0}
$$
\end{example}

\begin{lemma}
\label{lemma-rings-RS}
L'exemple \ref{example-rings}
satisfait la condition de Rim-Schlessinger (RS).
En particulier, $\Deformationcategory_P$ est une catégorie de déformations
pour toute $k$-algèbre $P$.
\end{lemma}

\begin{proof}
Soient $A_1 \to A$ et $A_2 \to A$ des morphismes de $\mathcal{C}_\Lambda$.
Supposons $A_2 \to A$ surjectif. D'après Théorie formelle
des déformations, lemme
\ref{formal-defos-lemma-RS-2-categorical}
il suffit de montrer que le foncteur
$\mathcal{F}(A_1 \times_A A_2) \to
\mathcal{F}(A_1) \times_{\mathcal{F}(A)} \mathcal{F}(A_2)$
est une équivalence de catégories.
C'est un cas particulier de Compléments sur l'algèbre, lemme
\ref{more-algebra-lemma-properties-algebras-over-fibre-product}.
\end{proof}

\begin{lemma}
\label{lemma-rings-TI}
Dans l'exemple \ref{example-rings}, soit $P$ une $k$-algèbre. Alors
$$
T\Deformationcategory_P = \text{Ext}^1_P(\NL_{P/k}, P)
\quad\text{et}\quad
\text{Inf}(\Deformationcategory_P) = \text{Der}_k(P, P)
$$
\end{lemma}

\begin{proof}
Rappelons que $\text{Inf}(\Deformationcategory_P)$ est l'ensemble des
automorphismes de la déformation triviale
$P[\epsilon] = P \otimes_k k[\epsilon]$ de $P$ sur $k[\epsilon]$
qui sont égaux à l'identité modulo $\epsilon$.
D'après Théorie des déformations, lemme \ref{defos-lemma-huge-diagram},
cet ensemble est égal à $\Hom_P(\Omega_{P/k}, P)$, qui est à son tour
égal à $\text{Der}_k(P, P)$ d'après
Algèbre, lemme \ref{algebra-lemma-universal-omega}.

\medskip\noindent
Rappelons que $T\Deformationcategory_P$ est l'ensemble des classes d'isomorphisme
des déformations plates $Q$ de $P$ sur $k[\epsilon]$, plus précisément
l'ensemble des classes d'isomorphisme de $\Deformationcategory_P(k[\epsilon])$.
Rappelons qu'une $k[\epsilon]$-algèbre $Q$ telle que $Q/\epsilon Q = P$
est plate sur $k[\epsilon]$ si et seulement si
$$
0 \to P \xrightarrow{\epsilon} Q \to P \to 0
$$
est exacte. Cela est démontré dans Compléments sur les morphismes, lemme
\ref{more-morphisms-lemma-deform}, et, plus généralement, dans
Théorie des déformations, lemme \ref{defos-lemma-deform-module}.
Nous pouvons donc appliquer
Théorie des déformations, lemme \ref{defos-lemma-choices},
pour voir que l'ensemble des classes d'isomorphisme de telles
déformations est égal à $\text{Ext}^1_P(\NL_{P/k}, P)$.
\end{proof}

\begin{lemma}
\label{lemma-smooth}
Dans l'exemple \ref{example-rings}, soit $P$ une $k$-algèbre lisse. Alors
$T\Deformationcategory_P = (0)$.
\end{lemma}

\begin{proof}
D'après le lemme \ref{lemma-rings-TI}, il faut montrer que
$\text{Ext}^1_P(\NL_{P/k}, P) = (0)$.
Puisque $k \to P$ est lisse, $\NL_{P/k}$ est quasi-isomorphe au
complexe constitué d'un
$P$-module projectif de type fini placé en degré $0$.
\end{proof}

\begin{lemma}
\label{lemma-finite-type-rings-TI}
Dans le lemme \ref{lemma-rings-TI}, si $P$ est une $k$-algèbre de type fini, alors
\begin{enumerate}
\item $\text{Inf}(\Deformationcategory_P)$ est de dimension finie si et seulement si
$\dim(P) = 0$ ;
\item $T\Deformationcategory_P$ est de dimension finie si
$\Spec(P) \to \Spec(k)$ est lisse sauf en un nombre fini de points.
\end{enumerate}
\end{lemma}

\begin{proof}
Preuve de (1). Considérons $\text{Der}_k(P, P)$ comme un $P$-module.
S'il est de dimension finie sur $k$, il est alors de longueur finie
comme $P$-module ; son support est donc contenu dans un nombre fini de
points fermés de $\Spec(P)$
(Algèbre, lemme \ref{algebra-lemma-simple-pieces}).
Puisque $\text{Der}_k(P, P) = \Hom_P(\Omega_{P/k}, P)$,
nous voyons que
$\text{Der}_k(P, P)_\mathfrak p = \text{Der}_k(P_\mathfrak p, P_\mathfrak p)$
pour tout idéal premier $\mathfrak p \subset P$
(cela utilise Algèbre, lemmes
\ref{algebra-lemma-differentials-localize},
\ref{algebra-lemma-differentials-finitely-presented}, et
\ref{algebra-lemma-hom-from-finitely-presented}).
Soit $\mathfrak p$ un idéal premier minimal de $P$
correspondant à une composante irréductible de dimension $d > 0$.
Alors $P_\mathfrak p$ est un anneau local artinien
essentiellement de type fini sur $k$, avec corps résiduel,
et $\Omega_{P_\mathfrak p/k}$ est non nul, par exemple d'après
Algèbre, lemme \ref{algebra-lemma-characterize-smooth-over-field}.
Tout module fini non nul sur un anneau local artinien
possède à la fois un sous-module et un module quotient isomorphes au corps résiduel.
Nous trouvons ainsi que
$\text{Der}_k(P_\mathfrak p, P_\mathfrak p) =
\Hom_{P_\mathfrak p}(\Omega_{P_\mathfrak p/k}, P_\mathfrak p)$
est lui aussi non nul. En combinant tout ce qui précède, nous obtenons (1).

\medskip\noindent
Preuve de (2). Pour un idéal premier $\mathfrak p$ de $P$, nous utiliserons l'égalité
$\NL_{P_\mathfrak p/k} = (\NL_{P/k})_\mathfrak p$
(Algèbre, lemme \ref{algebra-lemma-localize-NL})
ainsi que
l'égalité
$\text{Ext}_P^1(\NL_{P/k}, P)_\mathfrak p =
\text{Ext}_{P_\mathfrak p}^1(\NL_{P_\mathfrak p/k}, P_\mathfrak p)$
(Compléments sur l'algèbre, lemme
\ref{more-algebra-lemma-pseudo-coherence-and-base-change-ext}).
Étant donné un idéal premier $\mathfrak p \subset P$,
le morphisme $k \to P$ est lisse en $\mathfrak p$ si et seulement si
$(\NL_{P/k})_\mathfrak p$ est quasi-isomorphe
à un module projectif de type fini placé en degré $0$ (cela résulte
immédiatement de la définition d'un homomorphisme d'anneaux lisse, mais aussi
du résultat plus fort d'Algèbre, lemme \ref{algebra-lemma-smooth-at-point}).

\medskip\noindent
Supposons que $P$ soit lisse sur $k$ en dehors d'un nombre fini d'idéaux premiers.
Ces idéaux premiers « mauvais » sont alors les idéaux maximaux
$\mathfrak m_1, \ldots, \mathfrak m_n \subset P$, d'après
Algèbre, lemme \ref{algebra-lemma-finite-type-algebra-finite-nr-primes},
et parce que les idéaux premiers « mauvais » forment un fermé de $\Spec(P)$.
Pour $\mathfrak p \not \in \{\mathfrak m_1, \ldots, \mathfrak m_n\}$,
nous avons $\text{Ext}^1_P(\NL_{P/k}, P)_\mathfrak p = 0$ d'après ce qui précède.
Ainsi, $\text{Ext}^1_P(\NL_{P/k}, P)$ est un $P$-module fini
dont le support est contenu dans $\{\mathfrak m_1, \ldots, \mathfrak m_r\}$.
D'après Algèbre, proposition
\ref{algebra-proposition-minimal-primes-associated-primes}
par exemple, nous trouvons que la dimension sur $k$ de
$\text{Ext}^1_P(\NL_{P/k}, P)$ est une combinaison linéaire entière finie
des $\dim_k \kappa(\mathfrak m_i)$ ; elle est donc finie d'après
le théorème des zéros de Hilbert
(Algèbre, théorème \ref{algebra-theorem-nullstellensatz}).
\end{proof}

\noindent
Dans l'exemple \ref{example-rings}, soit $P$ une
$k$-algèbre de type fini. Alors $\Deformationcategory_P$
admet une présentation par un groupoïde lisse proreprésentable en foncteurs
sur $\mathcal{C}_\Lambda$ si et seulement si $\dim(P) = 0$.
De plus, $\Deformationcategory_P$ possède un objet
formel versel si $\Spec(P) \to \Spec(k)$ n'a qu'un nombre fini de
points singuliers. Cela résulte des lemmes \ref{lemma-rings-RS} et
\ref{lemma-finite-type-rings-TI},
ainsi que de la discussion générale de la section \ref{section-general}.

\begin{lemma}
\label{lemma-rings-hull}
Dans l'exemple \ref{example-rings}, supposons que $P$ soit une
$k$-algèbre de type fini telle que $\Spec(P) \to \Spec(k)$ soit lisse sauf
en un nombre fini de points.
Supposons que $\Lambda$ soit un anneau local complet de corps résiduel $k$
(cas classique). Alors le foncteur
$$
F : \mathcal{C}_\Lambda \longrightarrow \textit{Ensembles},\quad
A \longmapsto \Ob(\Deformationcategory_P(A))/\cong
$$
des classes d'isomorphisme d'objets possède une enveloppe.
\end{lemma}

\begin{proof}
Cela résulte immédiatement des lemmes \ref{lemma-rings-RS} et
\ref{lemma-finite-type-rings-TI}, ainsi que de
Théorie formelle des déformations, lemme \ref{formal-defos-lemma-RS-implies-S1-S2}
et remarque \ref{formal-defos-remark-compose-minimal-into-iso-classes}.
\end{proof}

\begin{lemma}
\label{lemma-localization}
Dans l'exemple \ref{example-rings}, soit $P$ une $k$-algèbre.
Soit $S \subset P$ une partie multiplicative. Il existe un foncteur naturel
$$
\Deformationcategory_P \longrightarrow \Deformationcategory_{S^{-1}P}
$$
de catégories de déformations.
\end{lemma}

\begin{proof}
Étant donnée une déformation de $P$, nous pouvons la localiser
pour obtenir une déformation de la localisation ; cela est
clair, et nous invitons le lecteur à omettre la démonstration. Plus précisément,
soit $(A, Q) \to (k, P)$ un morphisme de $\mathcal{F}$, c'est-à-dire
un objet de $\Deformationcategory_P$. Soit $S_Q \subset Q$ l'image
réciproque de $S$. Alors
Ainsi, $(A, S_Q^{-1}Q) \to (k, S^{-1}P)$
est l'objet recherché de $\Deformationcategory_{S^{-1}P}$.
\end{proof}

\begin{lemma}
\label{lemma-henselization}
Dans l'exemple \ref{example-rings}, soit $P$ une $k$-algèbre.
Soit $J \subset P$ un idéal.
Notons $(P^h, J^h)$ l'hensélisation du couple $(P, J)$.
Il existe un foncteur naturel
$$
\Deformationcategory_P \longrightarrow \Deformationcategory_{P^h}
$$
de catégories de déformations.
\end{lemma}

\begin{proof}
Étant donnée une déformation de $P$, nous pouvons en prendre l'hensélisation
pour obtenir une déformation de l'hensélisé ; cela est
clair, et nous invitons le lecteur à omettre la démonstration. Plus précisément,
soit $(A, Q) \to (k, P)$ un morphisme de $\mathcal{F}$, c'est-à-dire
un objet de $\Deformationcategory_P$. Notons $J_Q \subset Q$ l'image
réciproque de $J$ dans $Q$. Soit $(Q^h, J_Q^h)$ l'hensélisation
du couple $(Q, J_Q)$. Rappelons que $Q \to Q^h$ est plat
(Compléments sur l'algèbre, lemme \ref{more-algebra-lemma-henselization-flat})
et que $Q^h$ est donc plat sur $A$.
D'après Compléments sur l'algèbre, lemme \ref{more-algebra-lemma-henselization-integral},
nous voyons que l'application $Q^h \to P^h$ induit un isomorphisme
$Q^h \otimes_A k = Q^h \otimes_Q P = P^h$.
Ainsi, $(A, Q^h) \to (k, P^h)$ est l'objet recherché de
$\Deformationcategory_{P^h}$.
\end{proof}

\begin{lemma}
\label{lemma-strict-henselization}
Dans l'exemple \ref{example-rings}, soit $P$ une $k$-algèbre.
Supposons que $P$ soit un anneau local et soit $P^{sh}$ un hensélisé strict de $P$.
Il existe un foncteur naturel
$$
\Deformationcategory_P \longrightarrow \Deformationcategory_{P^{sh}}
$$
de catégories de déformations.
\end{lemma}

\begin{proof}
Étant donnée une déformation de $P$, nous pouvons en prendre l'hensélisé strict
pour obtenir une déformation de l'hensélisé strict ; cela est
clair, et nous invitons le lecteur à omettre la démonstration. Plus précisément,
soit $(A, Q) \to (k, P)$ un morphisme de $\mathcal{F}$, c'est-à-dire
un objet de $\Deformationcategory_P$. Puisque le noyau de la surjection
$Q \to P$ est nilpotent, nous trouvons que $Q$ est un anneau local de
même corps résiduel que $P$. Soit $Q^{sh}$ l'hensélisé strict
de $Q$. Rappelons que $Q \to Q^{sh}$ est plat
(Compléments sur l'algèbre, lemme \ref{more-algebra-lemma-dumb-properties-henselization})
et que $Q^{sh}$ est donc plat sur $A$.
D'après Algèbre, lemme \ref{algebra-lemma-quotient-strict-henselization},
nous voyons que l'application $Q^{sh} \to P^{sh}$ induit un isomorphisme
$Q^{sh} \otimes_A k = Q^{sh} \otimes_Q P = P^{sh}$.
Ainsi, $(A, Q^{sh}) \to (k, P^{sh})$ est l'objet recherché de
$\Deformationcategory_{P^{sh}}$.
\end{proof}

\begin{lemma}
\label{lemma-completion}
Dans l'exemple \ref{example-rings}, soit $P$ une $k$-algèbre.
Supposons $P$ noethérien et soit $J \subset P$ un idéal.
Notons $P^\wedge$ le complété $J$-adique.
Il existe un foncteur naturel
$$
\Deformationcategory_P \longrightarrow \Deformationcategory_{P^\wedge}
$$
de catégories de déformations.
\end{lemma}

\begin{proof}
Étant donnée une déformation de $P$, nous pouvons en prendre le complété
pour obtenir une déformation du complété ; cela est
clair, et nous invitons le lecteur à omettre la démonstration. Plus précisément,
soit $(A, Q) \to (k, P)$ un morphisme de $\mathcal{F}$, c'est-à-dire
un objet de $\Deformationcategory_P$. Observons que $Q$ est un anneau
noethérien : le noyau de l'homomorphisme surjectif d'anneaux $Q \to P$ est
nilpotent et de type fini, et $P$ est noethérien ; appliquons
Algèbre, lemme \ref{algebra-lemma-completion-Noetherian}.
Notons $J_Q \subset Q$ l'image
réciproque de $J$ dans $Q$. Soit $Q^\wedge$ le complété $J_Q$-adique de $Q$.
Rappelons que $Q \to Q^\wedge$ est plat
(Algèbre, lemme \ref{algebra-lemma-completion-flat})
et que $Q^\wedge$ est donc plat sur $A$.
L'application induite $Q^\wedge \to P^\wedge$ induit un isomorphisme
$Q^\wedge \otimes_A k = Q^\wedge \otimes_Q P = P^\wedge$ d'après
Algèbre, lemme \ref{algebra-lemma-completion-tensor}, par exemple.
Ainsi, $(A, Q^\wedge) \to (k, P^\wedge)$
est l'objet recherché de $\Deformationcategory_{P^\wedge}$.
\end{proof}

\begin{lemma}
\label{lemma-power-series-rings-TI}
Dans le lemme \ref{lemma-rings-TI}, si $P = k[[x_1, \ldots, x_n]]/(f)$
pour un certain $f \in (x_1, \ldots, x_n)^2$ non nul, alors
\begin{enumerate}
\item $\text{Inf}(\Deformationcategory_P)$ est de dimension finie
si et seulement si $n = 1$ ;
\item $T\Deformationcategory_P$ est de dimension finie si
$$
\sqrt{(f, \partial f/\partial x_1, \ldots,  \partial f/\partial x_n)} =
(x_1, \ldots, x_n)
$$
\end{enumerate}
\end{lemma}

\begin{proof}
Preuve de (1). Considérons les dérivations $\partial/\partial x_i$ de
$k[[x_1, \ldots, x_n]]$ sur $k$. Posons $f_i = \partial f/\partial x_i$.
La dérivation
$$
\theta = \sum h_i \partial/\partial x_i
$$
de $k[[x_1, \ldots, x_n]]$
induit une dérivation de $P = k[[x_1, \ldots, x_n]]/(f)$
si et seulement si
$\sum h_i f_i \in (f)$. De plus, la dérivation induite de $P$
est nulle si et seulement si $h_i \in (f)$ pour $i = 1, \ldots, n$.
Nous obtenons ainsi
$$
\Ker((f_1, \ldots, f_n) : P^{\oplus n} \longrightarrow P) \subset
\text{Der}_k(P, P)
$$
Le membre de gauche n'est un $k$-espace vectoriel de dimension finie que si
$n = 1$ ; nous omettons la démonstration. Nous laissons aussi au lecteur le soin de voir
que le membre de droite est de dimension finie si $n = 1$.
Cela démontre (1).

\medskip\noindent
Preuve de (2). Soit $Q$ une déformation plate de $P$ sur $k[\epsilon]$,
comme dans la démonstration du lemme \ref{lemma-rings-TI}. Choisissons des relèvements $q_i \in Q$
de l'image de $x_i$ dans $P$. Alors $Q$ est un anneau local complet
dont l'idéal maximal est engendré par $q_1, \ldots, q_n$ et $\epsilon$
(nous omettons un court argument). Nous obtenons ainsi une surjection
$$
k[\epsilon][[x_1, \ldots, x_n]] \longrightarrow Q,\quad
x_i \longmapsto q_i
$$
Choisissons un élément de la forme
$f + \epsilon g \in k[\epsilon][[x_1, \ldots, x_n]]$
qui s'envoie sur zéro dans $Q$. Observons que $g$ est bien défini modulo $(f)$.
Puisque $Q$ est plat sur $k[\epsilon]$, nous obtenons
$$
Q = k[\epsilon][[x_1, \ldots, x_n]]/(f + \epsilon g)
$$
Enfin, modifier le choix de $q_i$ revient à
remplacer les coordonnées $x_i$ par $x_i + \epsilon h_i$
pour certains $h_i \in k[[x_1, \ldots, x_n]]$. Alors
$f + \epsilon g$ devient $f + \epsilon (g + \sum h_i f_i)$,
où $f_i = \partial f/\partial x_i$. Nous voyons donc que la
classe d'isomorphisme de la déformation $Q$ est déterminée
par un élément de
$$
k[[x_1, \ldots, x_n]]/
(f, \partial f/\partial x_1, \ldots,  \partial f/\partial x_n)
$$
Ce quotient est de dimension finie sur $k$ si et seulement si
son support est le point fermé de $k[[x_1, \ldots, x_n]]$,
si et seulement si
$\sqrt{(f, \partial f/\partial x_1, \ldots,  \partial f/\partial x_n)} =
(x_1, \ldots, x_n)$.
\end{proof}






\section{Schémas}
\label{section-schemes}

\noindent
La théorie des déformations des schémas.

\begin{example}[Schémas]
\label{example-schemes}
Soit $\mathcal{F}$ la catégorie définie comme suit :
\begin{enumerate}
\item un objet est un couple $(A, X)$ formé d'un
objet $A$ de $\mathcal{C}_\Lambda$ et d'un schéma $X$ plat sur $A$ ;
\item un morphisme $(f, g) : (B, Y) \to (A, X)$ est formé d'un
morphisme $f : B \to A$ dans $\mathcal{C}_\Lambda$ et
d'un morphisme $g : X \to Y$ tel que
$$
\xymatrix{
X \ar[r]_g \ar[d] & Y \ar[d] \\
\Spec(A) \ar[r]^f & \Spec(B)
}
$$
soit un diagramme commutatif cartésien de schémas.
\end{enumerate}
Le foncteur $p : \mathcal{F} \to \mathcal{C}_\Lambda$ envoie $(A, X)$
sur $A$ et $(f, g)$ sur $f$. Il est clair que $p$ est cofibré en groupoïdes.
Étant donné un schéma $X$ sur $k$, soit $x_0 = (k, X)$ l'objet correspondant
de $\mathcal{F}(k)$. On pose
$$
\Deformationcategory_X = \mathcal{F}_{x_0}
$$
\end{example}

\begin{lemma}
\label{lemma-schemes-RS}
L'exemple \ref{example-schemes}
satisfait la condition de Rim-Schlessinger (RS).
En particulier, $\Deformationcategory_X$ est une catégorie de déformations
pour tout schéma $X$ sur $k$.
\end{lemma}

\begin{proof}
Soient $A_1 \to A$ et $A_2 \to A$ des morphismes de $\mathcal{C}_\Lambda$.
Supposons $A_2 \to A$ surjectif. D'après Théorie formelle
des déformations, lemme
\ref{formal-defos-lemma-RS-2-categorical}
il suffit de montrer que le foncteur
$\mathcal{F}(A_1 \times_A A_2) \to
\mathcal{F}(A_1) \times_{\mathcal{F}(A)} \mathcal{F}(A_2)$
est une équivalence de catégories.
Observons que
$$
\xymatrix{
\Spec(A) \ar[r] \ar[d] & \Spec(A_2) \ar[d] \\
\Spec(A_1) \ar[r] &
\Spec(A_1 \times_A A_2)
}
$$
est un diagramme cocartésien comme dans Compléments sur les morphismes, lemme
\ref{more-morphisms-lemma-pushout-along-thickening}.
Le lemme est donc un cas particulier de Compléments sur les morphismes, lemme
\ref{more-morphisms-lemma-equivalence-categories-schemes-over-pushout-flat}.
\end{proof}

\begin{lemma}
\label{lemma-schemes-TI}
Dans l'exemple \ref{example-schemes}, soit $X$ un schéma sur $k$. Alors
$$
\text{Inf}(\Deformationcategory_X) =
\text{Ext}^0_{\mathcal{O}_X}(\NL_{X/k}, \mathcal{O}_X) =
\Hom_{\mathcal{O}_X}(\Omega_{X/k}, \mathcal{O}_X) =
\text{Der}_k(\mathcal{O}_X, \mathcal{O}_X)
$$
et
$$
T\Deformationcategory_X =
\text{Ext}^1_{\mathcal{O}_X}(\NL_{X/k}, \mathcal{O}_X)
$$
\end{lemma}

\begin{proof}
Rappelons que $\text{Inf}(\Deformationcategory_X)$ est l'ensemble des
automorphismes de la déformation triviale
$X' = X \times_{\Spec(k)} \Spec(k[\epsilon])$ de $X$ sur $k[\epsilon]$
qui sont égaux à l'identité modulo $\epsilon$.
D'après Théorie des déformations, lemme \ref{defos-lemma-deform},
cet ensemble est égal à $\text{Ext}^0_{\mathcal{O}_X}(\NL_{X/k}, \mathcal{O}_X)$.
L'égalité $\text{Ext}^0_{\mathcal{O}_X}(\NL_{X/k}, \mathcal{O}_X) =
\Hom_{\mathcal{O}_X}(\Omega_{X/k}, \mathcal{O}_X)$ résulte de
Compléments sur les morphismes, lemme
\ref{more-morphisms-lemma-netherlander-quasi-coherent}.
L'égalité
$\Hom_{\mathcal{O}_X}(\Omega_{X/k}, \mathcal{O}_X) =
\text{Der}_k(\mathcal{O}_X, \mathcal{O}_X)$
résulte de Morphismes, lemme
\ref{morphisms-lemma-universal-derivation-universal}.

\medskip\noindent
Rappelons que $T_{x_0}\Deformationcategory_X$ est l'ensemble des classes d'isomorphisme
des déformations plates $X'$ de $X$ sur $k[\epsilon]$, plus précisément
l'ensemble des classes d'isomorphisme de $\Deformationcategory_X(k[\epsilon])$.
La seconde assertion du lemme résulte donc de
Théorie des déformations, lemme \ref{defos-lemma-deform}.
\end{proof}

\begin{lemma}
\label{lemma-proper-schemes-TI}
Dans le lemme \ref{lemma-schemes-TI}, si $X$ est propre sur $k$, alors
$\text{Inf}(\Deformationcategory_X)$ et $T\Deformationcategory_X$ sont
de dimension finie.
\end{lemma}

\begin{proof}
D'après le lemme, il faut montrer que
$\Ext^1_{\mathcal{O}_X}(\NL_{X/k}, \mathcal{O}_X)$ et
$\Ext^0_{\mathcal{O}_X}(\NL_{X/k}, \mathcal{O}_X)$ sont de
dimension finie. D'après Compléments sur les morphismes, lemme
\ref{more-morphisms-lemma-netherlander-fp}
et le fait que $X$ est noethérien, nous voyons que
$\NL_{X/k}$ a des faisceaux de cohomologie cohérents, nuls sauf
en degrés $0$ et $-1$.
D'après Catégories dérivées de schémas, lemme \ref{perfect-lemma-ext-finite},
les groupes $\Ext$ affichés sont des $k$-espaces vectoriels de dimension finie,
ce qui achève la démonstration.
\end{proof}

\noindent
Dans l'exemple \ref{example-schemes}, si $X$ est un schéma propre sur $k$,
alors $\Deformationcategory_X$
admet une présentation par un groupoïde lisse proreprésentable en foncteurs
sur $\mathcal{C}_\Lambda$
et possède a fortiori un objet formel versel (minimal). Cela résulte
des lemmes \ref{lemma-schemes-RS} et
\ref{lemma-proper-schemes-TI},
ainsi que de la discussion générale de la section \ref{section-general}.

\begin{lemma}
\label{lemma-schemes-hull}
Dans l'exemple \ref{example-schemes}, supposons que $X$ soit un $k$-schéma propre.
Supposons que $\Lambda$ soit un anneau local complet de corps résiduel $k$
(cas classique). Alors le foncteur
$$
F : \mathcal{C}_\Lambda \longrightarrow \textit{Ensembles},\quad
A \longmapsto \Ob(\Deformationcategory_X(A))/\cong
$$
des classes d'isomorphisme d'objets possède une enveloppe. Si
$\text{Der}_k(\mathcal{O}_X, \mathcal{O}_X) = 0$, alors
$F$ est proreprésentable.
\end{lemma}

\begin{proof}
L'existence d'une enveloppe résulte immédiatement des
lemmes \ref{lemma-schemes-RS} et \ref{lemma-proper-schemes-TI}, ainsi que de
Théorie formelle des déformations, lemme \ref{formal-defos-lemma-RS-implies-S1-S2}
et remarque \ref{formal-defos-remark-compose-minimal-into-iso-classes}.

\medskip\noindent
Supposons $\text{Der}_k(\mathcal{O}_X, \mathcal{O}_X) = 0$. Alors
$\Deformationcategory_X$ et $F$ sont équivalents d'après
Théorie formelle des déformations, lemme \ref{formal-defos-lemma-infdef-trivial}.
Ainsi, $F$ est un foncteur de déformations (car $\Deformationcategory_X$ est une
catégorie de déformations) d'espace tangent de dimension finie, et nous pouvons appliquer
Théorie formelle des déformations, théorème
\ref{formal-defos-theorem-Schlessinger-prorepresentability}.
\end{proof}

\begin{lemma}
\label{lemma-open}
Dans l'exemple \ref{example-schemes}, soit $X$ un schéma sur $k$.
Soit $U \subset X$ un sous-schéma ouvert.
Il existe un foncteur naturel
$$
\Deformationcategory_X \longrightarrow \Deformationcategory_U
$$
de catégories de déformations.
\end{lemma}

\begin{proof}
Étant donnée une déformation de $X$, nous pouvons en prendre l'ouvert correspondant
pour obtenir une déformation de $U$. Nous omettons les détails.
\end{proof}

\begin{lemma}
\label{lemma-affine}
Dans l'exemple \ref{example-schemes}, soit $X = \Spec(P)$ un
schéma affine sur $k$. Avec $\Deformationcategory_P$ comme dans
l'exemple \ref{example-rings}, il existe une équivalence naturelle
$$
\Deformationcategory_X \longrightarrow \Deformationcategory_P
$$
de catégories de déformations.
\end{lemma}

\begin{proof}
Le foncteur envoie $(A, Y)$ sur $\Gamma(Y, \mathcal{O}_Y)$.
Cette construction convient parce que
toute déformation de $X$ est affine d'après
Compléments sur les morphismes, lemme \ref{more-morphisms-lemma-thickening-affine-scheme}.
\end{proof}

\begin{lemma}
\label{lemma-local-ring}
Dans l'exemple \ref{example-schemes}, soit $X$ un schéma sur $k$.
Soit $p \in X$ un point. Avec $\Deformationcategory_{\mathcal{O}_{X, p}}$
comme dans l'exemple \ref{example-rings}, il existe un foncteur naturel
$$
\Deformationcategory_X
\longrightarrow
\Deformationcategory_{\mathcal{O}_{X, p}}
$$
de catégories de déformations.
\end{lemma}

\begin{proof}
Choisissons un ouvert affine $U = \Spec(P) \subset X$ contenant $p$.
Alors $\mathcal{O}_{X, p}$ est une localisation de $P$. Composons
les foncteurs des
lemmes \ref{lemma-open}, \ref{lemma-affine} et \ref{lemma-localization}.
\end{proof}

\begin{situation}
\label{situation-glueing}
Soit $\Lambda \to k$ comme dans la section \ref{section-general}.
Soit $X$ un schéma sur $k$ qui possède un recouvrement ouvert affine
$X = U_1 \cup U_2$, où $U_{12} = U_1 \cap U_2$ est lui aussi affine.
Écrivons $U_1 = \Spec(P_1)$, $U_2 = \Spec(P_2)$ et $U_{12} = \Spec(P_{12})$.
Soient $\Deformationcategory_X$, $\Deformationcategory_{U_1}$,
$\Deformationcategory_{U_2}$ et $\Deformationcategory_{U_{12}}$
comme dans l'exemple \ref{example-schemes}, et soient
$\Deformationcategory_{P_1}$, $\Deformationcategory_{P_2}$ et
$\Deformationcategory_{P_{12}}$ comme dans l'exemple \ref{example-rings}.
\end{situation}

\begin{lemma}
\label{lemma-glueing}
Dans la situation \ref{situation-glueing},
il existe une équivalence
$$
\Deformationcategory_X =
\Deformationcategory_{P_1}
\times_{\Deformationcategory_{P_{12}}}
\Deformationcategory_{P_2}
$$
de catégories de déformations ; voir les exemples \ref{example-schemes} et
\ref{example-rings}.
\end{lemma}

\begin{proof}
Il suffit de montrer que les foncteurs du lemme \ref{lemma-open}
définissent une équivalence
$$
\Deformationcategory_X \longrightarrow
\Deformationcategory_{U_1}
\times_{\Deformationcategory_{U_{12}}}
\Deformationcategory_{U_2}
$$
car nous pouvons alors appliquer le lemme \ref{lemma-affine} pour passer aux anneaux.
Pour cela, construisons un quasi-inverse. Notons
$F_i : \Deformationcategory_{U_i} \to \Deformationcategory_{U_{12}}$
le foncteur du lemme \ref{lemma-open}.
Un objet du membre de droite est donné par un $A$ dans $\mathcal{C}_\Lambda$,
des objets $(A, V_1) \to (k, U_1)$ et $(A, V_2) \to (k, U_2)$, et
un morphisme
$$
g : F_1(A, V_1) \to F_2(A, V_2)
$$
Ici, $F_i(A, V_i) = (A, V_{i, 3 - i})$, où $V_{i, 3 - i} \subset V_i$
est le sous-schéma ouvert dont le changement de base à $k$ est $U_{12} \subset U_i$.
Le morphisme $g$ définit un isomorphisme
$V_{1, 2} \to V_{2, 1}$ de schémas sur $A$, compatible
avec $\text{id} : U_{12} \to U_{12}$ sur $k$.
Ainsi, $(\{1, 2\}, V_i, V_{i, 3 - i}, g, g^{-1})$ est une donnée
de recollement comme dans Schémas, section \ref{schemes-section-glueing-schemes}.
Soit $Y$ le schéma recollé ; voir Schémas, lemme \ref{schemes-lemma-glue}.
Alors $Y$ est un schéma sur $A$, et les
compatibilités mentionnées ci-dessus montrent
qu'il existe un isomorphisme canonique
$Y \times_{\Spec(A)} \Spec(k) = X$.
Ainsi, $(A, Y) \to (k, X)$ est un objet de $\Deformationcategory_X$.
Nous omettons de vérifier que cette construction définit un foncteur
quasi-inverse de celui qui était donné.
\end{proof}






\section{Morphismes de schémas}
\label{section-schemes-morphisms}

\noindent
La théorie des déformations des morphismes de schémas.
Il ne s'agit bien sûr que d'un exemple de
déformations de diagrammes de schémas.

\begin{example}[Morphismes de schémas]
\label{example-schemes-morphisms}
Soit $\mathcal{F}$ la catégorie définie comme suit :
\begin{enumerate}
\item un objet est un couple $(A, X \to Y)$ formé d'un
objet $A$ de $\mathcal{C}_\Lambda$ et d'un morphisme
$X \to Y$ de schémas sur $A$, où $X$ et $Y$ sont tous deux plats sur $A$ ;
\item un morphisme $(f, g, h) : (A', X' \to Y') \to (A, X \to Y)$ est formé d'un
morphisme $f : A' \to A$ dans $\mathcal{C}_\Lambda$ et
de morphismes de schémas $g : X \to X'$ et $h : Y \to Y'$ tels que
$$
\xymatrix{
X \ar[r]_g \ar[d] & X' \ar[d] \\
Y \ar[r]_h \ar[d] & Y' \ar[d] \\
\Spec(A) \ar[r]^f & \Spec(A')
}
$$
soit un diagramme commutatif de schémas dont les deux carrés sont cartésiens.
\end{enumerate}
Le foncteur $p : \mathcal{F} \to \mathcal{C}_\Lambda$ envoie $(A, X \to Y)$
sur $A$ et $(f, g, h)$ sur $f$. Il est clair que $p$ est cofibré en groupoïdes.
Étant donné un morphisme de schémas $X \to Y$ sur $k$, soit $x_0 = (k, X \to Y)$
l'objet correspondant de $\mathcal{F}(k)$. On pose
$$
\Deformationcategory_{X \to Y} = \mathcal{F}_{x_0}
$$
\end{example}

\begin{lemma}
\label{lemma-schemes-morphisms-RS}
L'exemple \ref{example-schemes-morphisms}
satisfait la condition de Rim-Schlessinger (RS).
En particulier, $\Deformationcategory_{X \to Y}$ est une catégorie de déformations
pour tout morphisme de schémas $X \to Y$ sur $k$.
\end{lemma}

\begin{proof}
Soient $A_1 \to A$ et $A_2 \to A$ des morphismes de $\mathcal{C}_\Lambda$.
Supposons $A_2 \to A$ surjectif. D'après Théorie formelle
des déformations, lemme
\ref{formal-defos-lemma-RS-2-categorical}
il suffit de montrer que le foncteur
$\mathcal{F}(A_1 \times_A A_2) \to
\mathcal{F}(A_1) \times_{\mathcal{F}(A)} \mathcal{F}(A_2)$
est une équivalence de catégories.
Observons que
$$
\xymatrix{
\Spec(A) \ar[r] \ar[d] & \Spec(A_2) \ar[d] \\
\Spec(A_1) \ar[r] &
\Spec(A_1 \times_A A_2)
}
$$
est un diagramme cocartésien comme dans Compléments sur les morphismes, lemme
\ref{more-morphisms-lemma-pushout-along-thickening}.
Le lemme résulte donc immédiatement de
Compléments sur les morphismes, lemme
\ref{more-morphisms-lemma-equivalence-categories-schemes-over-pushout-flat}
car ce résultat décrit la catégorie des schémas plats sur $A_1 \times_A A_2$
comme le produit fibré de la catégorie des schémas plats sur $A_1$
et de la catégorie des schémas plats sur $A_2$ au-dessus de la catégorie des
schémas plats sur $A$.
\end{proof}

\begin{lemma}
\label{lemma-schemes-morphisms-TI}
Dans l'exemple \ref{example-schemes}, soit $f : X \to Y$ un morphisme de schémas
sur $k$. Il existe une suite exacte canonique de $k$-espaces vectoriels
$$
\xymatrix{
0 \ar[r] &
\text{Inf}(\Deformationcategory_{X \to Y}) \ar[r] &
\text{Inf}(\Deformationcategory_X \times \Deformationcategory_Y) \ar[r] &
\text{Der}_k(\mathcal{O}_Y, f_*\mathcal{O}_X) \ar[lld] \\
& T\Deformationcategory_{X \to Y} \ar[r] &
T(\Deformationcategory_X \times \Deformationcategory_Y) \ar[r] &
\text{Ext}^1_{\mathcal{O}_X}(Lf^*\NL_{Y/k}, \mathcal{O}_X)
}
$$
\end{lemma}

\begin{proof}
Le morphisme évident de catégories de déformations
$\Deformationcategory_{X \to Y} \to
\Deformationcategory_X \times \Deformationcategory_Y$
donne deux des flèches de la suite exacte du lemme.
Rappelons que $\text{Inf}(\Deformationcategory_{X \to Y})$
est l'ensemble des automorphismes de la déformation triviale
$$
f' :  X' = X \times_{\Spec(k)} \Spec(k[\epsilon])
\xrightarrow{f \times \text{id}}
Y' = Y \times_{\Spec(k)} \Spec(k[\epsilon])
$$
de $X \to Y$ sur $k[\epsilon]$ qui sont égaux à l'identité modulo $\epsilon$.
Il s'agit clairement des couples
$(\alpha, \beta) \in
\text{Inf}(\Deformationcategory_X \times \Deformationcategory_Y)$
d'automorphismes infinitésimaux de $X$ et $Y$ compatibles avec $f'$, c'est-à-dire
tels que $f' \circ \alpha = \beta \circ f'$.
D'après Théorie des déformations, lemme \ref{defos-lemma-huge-diagram-ringed-spaces},
pour un couple arbitraire $(\alpha, \beta)$, la différence entre
le morphisme $f' : X' \to Y'$ et le morphisme
$\beta^{-1} \circ f' \circ \alpha : X' \to Y'$ définit un élément
de
$$
\text{Der}_k(\mathcal{O}_Y, f_*\mathcal{O}_X) =
\Hom_{\mathcal{O}_Y}(\Omega_{Y/k}, f_*\mathcal{O}_X)
$$
L'égalité résulte de Compléments sur les morphismes, lemme
\ref{more-morphisms-lemma-netherlander-quasi-coherent}.
Cela définit la dernière flèche horizontale supérieure et montre l'exactitude
aux deux premières places. Pour l'application
$$
\text{Der}_k(\mathcal{O}_Y, f_*\mathcal{O}_X)
\to
T\Deformationcategory_{X \to Y}
$$
interprétons les éléments de la source comme des morphismes
$f_\epsilon : X' \to Y'$ sur $\Spec(k[\epsilon])$
égaux à $f$ modulo $\epsilon$,
à l'aide de Théorie des déformations, lemme \ref{defos-lemma-huge-diagram-ringed-spaces}.
Nous envoyons $f_\epsilon$ sur la classe d'isomorphisme de
$(f_\epsilon : X' \to Y')$ dans $T\Deformationcategory_{X \to Y}$.
Notons que $(f_\epsilon : X' \to Y')$ est isomorphe à la
déformation triviale $(f' : X' \to Y')$ exactement lorsque
$f_\epsilon  = \beta^{-1} \circ f \circ \alpha$ pour un certain
couple $(\alpha, \beta)$, ce qui implique l'exactitude à la troisième place.
Clairement, si une déformation du premier ordre
$(f_\epsilon : X_\epsilon \to Y_\epsilon)$
s'envoie sur zéro dans $T(\Deformationcategory_X \times \Deformationcategory_Y)$,
nous pouvons alors choisir des isomorphismes $X' \to X_\epsilon$ et $Y' \to Y_\epsilon$
et nous concluons que nous sommes dans l'image de la flèche sud-ouest.
Nous avons donc l'exactitude à la quatrième place.
Enfin, étant données deux déformations du premier ordre $X_\epsilon$, $Y_\epsilon$
de $X$, $Y$, il existe une obstruction dans
$$
ob(X_\epsilon, Y_\epsilon) \in
\text{Ext}^1_{\mathcal{O}_X}(Lf^*\NL_{Y/k}, \mathcal{O}_X)
$$
qui s'annule si et seulement si $f : X \to Y$ se relève en
$X_\epsilon \to Y_\epsilon$ ; voir
Théorie des déformations, lemme \ref{defos-lemma-huge-diagram-ringed-spaces}.
Cela achève la démonstration.
\end{proof}

\begin{lemma}
\label{lemma-proper-schemes-morphisms-TI}
Dans le lemme \ref{lemma-schemes-morphisms-TI}, si $X$ et $Y$ sont tous deux
propres sur $k$, alors
$\text{Inf}(\Deformationcategory_{X \to Y})$ et
$T\Deformationcategory_{X \to Y}$ sont de dimension finie.
\end{lemma}

\begin{proof}
Démonstration omise. Indication : raisonner comme dans le lemme \ref{lemma-proper-schemes-TI}
et utiliser la suite exacte du lemme.
\end{proof}

\noindent
Dans l'exemple \ref{example-schemes-morphisms},
si $X \to Y$ est un morphisme de schémas propres sur $k$,
alors $\Deformationcategory_{X \to Y}$
admet une présentation par un groupoïde lisse proreprésentable en foncteurs
sur $\mathcal{C}_\Lambda$
et possède a fortiori un objet formel versel (minimal). Cela résulte
des lemmes \ref{lemma-schemes-morphisms-RS} et
\ref{lemma-proper-schemes-morphisms-TI},
ainsi que de la discussion générale de la section \ref{section-general}.

\begin{lemma}
\label{lemma-schemes-morphisms-hull}
Dans l'exemple \ref{example-schemes-morphisms}, supposons que $X \to Y$
soit un morphisme de $k$-schémas propres.
Supposons que $\Lambda$ soit un anneau local complet de corps résiduel $k$
(cas classique). Alors le foncteur
$$
F : \mathcal{C}_\Lambda \longrightarrow \textit{Ensembles},\quad
A \longmapsto \Ob(\Deformationcategory_{X \to Y}(A))/\cong
$$
des classes d'isomorphisme d'objets possède une enveloppe. Si
$\text{Der}_k(\mathcal{O}_X, \mathcal{O}_X) =
\text{Der}_k(\mathcal{O}_Y, \mathcal{O}_Y) = 0$, alors
$F$ est proreprésentable.
\end{lemma}

\begin{proof}
L'existence d'une enveloppe résulte immédiatement des
lemmes \ref{lemma-schemes-morphisms-RS} et
\ref{lemma-proper-schemes-morphisms-TI}, ainsi que de
Théorie formelle des déformations, lemme \ref{formal-defos-lemma-RS-implies-S1-S2}
et remarque \ref{formal-defos-remark-compose-minimal-into-iso-classes}.

\medskip\noindent
Supposons $\text{Der}_k(\mathcal{O}_X, \mathcal{O}_X) =
\text{Der}_k(\mathcal{O}_Y, \mathcal{O}_Y) = 0$. Alors
la suite exacte du lemme \ref{lemma-schemes-morphisms-TI},
combinée au lemme \ref{lemma-schemes-TI},
montre que $\text{Inf}(\Deformationcategory_{X \to Y}) = 0$.
Alors $\Deformationcategory_{X \to Y}$ et $F$ sont équivalents d'après
Théorie formelle des déformations, lemme \ref{formal-defos-lemma-infdef-trivial}.
Ainsi, $F$ est un foncteur de déformations (car
$\Deformationcategory_{X \to Y}$ est une
catégorie de déformations) d'espace tangent de dimension finie, et nous pouvons appliquer
Théorie formelle des déformations, théorème
\ref{formal-defos-theorem-Schlessinger-prorepresentability}.
\end{proof}

\begin{lemma}
\label{lemma-schemes-morphisms-smooth-to-base}
\begin{reference}
Cette question est traitée dans \cite[Section 5.3]{Ravi-Murphys-Law} et
\cite[Theorem 3.3]{Ran-deformations}.
\end{reference}
Dans l'exemple \ref{example-schemes}, soit $f : X \to Y$ un morphisme de schémas
sur $k$. Si $f_*\mathcal{O}_X = \mathcal{O}_Y$ et $R^1f_*\mathcal{O}_X = 0$,
alors le morphisme de catégories de déformations
$$
\Deformationcategory_{X \to Y} \to \Deformationcategory_X
$$
est une équivalence.
\end{lemma}

\begin{proof}
Construisons un quasi-inverse du foncteur d'oubli du lemme.
Plus précisément, supposons que $(A, U)$ soit un objet de $\Deformationcategory_X$.
L'application donnée $X \to U$ est un épaississement d'ordre fini, que nous pouvons utiliser
pour identifier les espaces topologiques sous-jacents à $U$ et $X$ ; voir
Compléments sur les morphismes, section \ref{more-morphisms-section-thickenings}.
Nous pouvons donc considérer, et considérerons, $\mathcal{O}_U$ comme un faisceau de
$A$-algèbres sur $X$ ; de plus, la platitude de $U \to \Spec(A)$
signifie que $\mathcal{O}_U$ est plat comme faisceau de $A$-modules.
En particulier, nous avons une filtration
$$
0 = \mathfrak m_A^n\mathcal{O}_U \subset
\mathfrak m_A^{n - 1}\mathcal{O}_U \subset \ldots \subset
\mathfrak m_A^2\mathcal{O}_U \subset
\mathfrak m_A\mathcal{O}_U \subset \mathcal{O}_U
$$
dont les sous-quotients sont égaux à
$\mathcal{O}_X \otimes_k \mathfrak m_A^i/\mathfrak m_A^{i + 1}$
par platitude ; voir Compléments sur les morphismes, lemme \ref{more-morphisms-lemma-deform},
ou, plus généralement, Théorie des déformations, lemme \ref{defos-lemma-deform-module}.
Posons
$$
\mathcal{O}_V = f_*\mathcal{O}_U
$$
considéré comme un faisceau de $A$-algèbres sur $Y$. Puisque
$R^1f_*\mathcal{O}_X = 0$, la description ci-dessus montre que
$R^1f_*(\mathfrak m_A^i\mathcal{O}_U/\mathfrak m_A^{i + 1}\mathcal{O}_U) = 0$
pour tout $i$. Il s'ensuit que les suites
$$
0 \to
(f_*\mathcal{O}_X) \otimes_k \mathfrak m_A^i/\mathfrak m_A^{i + 1} \to
f_*(\mathcal{O}_U/\mathfrak m_A^{i + 1}\mathcal{O}_U) \to
f_*(\mathcal{O}_U/\mathfrak m_A^i\mathcal{O}_U) \to 0
$$
sont exactes pour tout $i$. En parcourant en sens inverse les références données ci-dessus
(et en raisonnant par récurrence), nous trouvons que $\mathcal{O}_V$ est un faisceau plat
de $A$-algèbres tel que
$\mathcal{O}_V/\mathfrak m_A\mathcal{O}_V = \mathcal{O}_Y$.
À l'aide de Compléments sur les morphismes, lemme
\ref{more-morphisms-lemma-first-order-thickening}
nous trouvons que $(Y, \mathcal{O}_V)$ est un schéma, que nous notons $V$.
L'égalité $\mathcal{O}_V = f_*\mathcal{O}_U$ définit un
morphisme d'espaces annelés $U \to V$ dont on voit aisément qu'il est
un morphisme de schémas. La platitude déjà établie
achève la démonstration.
\end{proof}







\section{Espaces algébriques}
\label{section-algebraic-spaces}

\noindent
La théorie des déformations des espaces algébriques.

\begin{example}[Espaces algébriques]
\label{example-spaces}
Soit $\mathcal{F}$ la catégorie définie comme suit :
\begin{enumerate}
\item un objet est un couple $(A, X)$ formé d'un
objet $A$ de $\mathcal{C}_\Lambda$ et d'un espace algébrique
$X$ plat sur $A$ ;
\item un morphisme $(f, g) : (B, Y) \to (A, X)$ est formé d'un
morphisme $f : B \to A$ dans $\mathcal{C}_\Lambda$ et
d'un morphisme $g : X \to Y$ d'espaces algébriques sur $\Lambda$
tel que
$$
\xymatrix{
X \ar[r]_g \ar[d] & Y \ar[d] \\
\Spec(A) \ar[r]^f & \Spec(B)
}
$$
soit un diagramme commutatif cartésien d'espaces algébriques.
\end{enumerate}
Le foncteur $p : \mathcal{F} \to \mathcal{C}_\Lambda$ envoie $(A, X)$
sur $A$ et $(f, g)$ sur $f$. Il est clair que $p$ est cofibré en groupoïdes.
Étant donné un espace algébrique $X$ sur $k$, soit
$x_0 = (k, X)$ l'objet correspondant de $\mathcal{F}(k)$. On pose
$$
\Deformationcategory_X = \mathcal{F}_{x_0}
$$
\end{example}

\begin{lemma}
\label{lemma-spaces-RS}
L'exemple \ref{example-spaces}
satisfait la condition de Rim-Schlessinger (RS).
En particulier, $\Deformationcategory_X$ est une catégorie de déformations
pour tout espace algébrique $X$ sur $k$.
\end{lemma}

\begin{proof}
Soient $A_1 \to A$ et $A_2 \to A$ des morphismes de $\mathcal{C}_\Lambda$.
Supposons $A_2 \to A$ surjectif. D'après Théorie formelle
des déformations, lemme
\ref{formal-defos-lemma-RS-2-categorical}
il suffit de montrer que le foncteur
$\mathcal{F}(A_1 \times_A A_2) \to
\mathcal{F}(A_1) \times_{\mathcal{F}(A)} \mathcal{F}(A_2)$
est une équivalence de catégories.
Observons que
$$
\xymatrix{
\Spec(A) \ar[r] \ar[d] & \Spec(A_2) \ar[d] \\
\Spec(A_1) \ar[r] &
\Spec(A_1 \times_A A_2)
}
$$
est un diagramme cocartésien comme dans Sommes amalgamées d'espaces, lemme
\ref{spaces-pushouts-lemma-pushout-along-thickening}.
Le lemme est donc un cas particulier de Sommes amalgamées d'espaces, lemme
\ref{spaces-pushouts-lemma-equivalence-categories-spaces-pushout-flat}.
\end{proof}

\begin{lemma}
\label{lemma-spaces-TI}
Dans l'exemple \ref{example-spaces}, soit $X$ un espace algébrique sur $k$. Alors
$$
\text{Inf}(\Deformationcategory_X) =
\text{Ext}^0_{\mathcal{O}_X}(\NL_{X/k}, \mathcal{O}_X) =
\Hom_{\mathcal{O}_X}(\Omega_{X/k}, \mathcal{O}_X) =
\text{Der}_k(\mathcal{O}_X, \mathcal{O}_X)
$$
et
$$
T\Deformationcategory_X =
\text{Ext}^1_{\mathcal{O}_X}(\NL_{X/k}, \mathcal{O}_X)
$$
\end{lemma}

\begin{proof}
Rappelons que $\text{Inf}(\Deformationcategory_X)$ est l'ensemble des
automorphismes de la déformation triviale
$X' = X \times_{\Spec(k)} \Spec(k[\epsilon])$ de $X$ sur $k[\epsilon]$
qui sont égaux à l'identité modulo $\epsilon$.
D'après Théorie des déformations, lemme \ref{defos-lemma-deform-spaces},
cet ensemble est égal à $\text{Ext}^0_{\mathcal{O}_X}(\NL_{X/k}, \mathcal{O}_X)$.
L'égalité $\text{Ext}^0_{\mathcal{O}_X}(\NL_{X/k}, \mathcal{O}_X) =
\Hom_{\mathcal{O}_X}(\Omega_{X/k}, \mathcal{O}_X)$ résulte de
Compléments sur les morphismes d'espaces, lemme
\ref{spaces-more-morphisms-lemma-netherlander-quasi-coherent}.
L'égalité
$\Hom_{\mathcal{O}_X}(\Omega_{X/k}, \mathcal{O}_X) =
\text{Der}_k(\mathcal{O}_X, \mathcal{O}_X)$
résulte de Compléments sur les morphismes d'espaces, définition
\ref{spaces-more-morphisms-definition-sheaf-differentials}, et de
Modules sur les sites, définition
\ref{sites-modules-definition-module-differentials}.

\medskip\noindent
Rappelons que $T_{x_0}\Deformationcategory_X$ est l'ensemble des classes d'isomorphisme
des déformations plates $X'$ de $X$ sur $k[\epsilon]$, plus précisément
l'ensemble des classes d'isomorphisme de $\Deformationcategory_X(k[\epsilon])$.
La seconde assertion du lemme résulte donc de
Théorie des déformations, lemme \ref{defos-lemma-deform-spaces}.
\end{proof}

\begin{lemma}
\label{lemma-proper-spaces-TI}
Dans le lemme \ref{lemma-spaces-TI}, si $X$ est propre sur $k$, alors
$\text{Inf}(\Deformationcategory_X)$ et $T\Deformationcategory_X$ sont
de dimension finie.
\end{lemma}

\begin{proof}
D'après le lemme, il faut montrer que
$\Ext^1_{\mathcal{O}_X}(\NL_{X/k}, \mathcal{O}_X)$ et
$\Ext^0_{\mathcal{O}_X}(\NL_{X/k}, \mathcal{O}_X)$ sont de
dimension finie. D'après Compléments sur les morphismes d'espaces, lemme
\ref{spaces-more-morphisms-lemma-netherlander-fp}
et le fait que $X$ est noethérien, nous voyons que
$\NL_{X/k}$ a des faisceaux de cohomologie cohérents, nuls sauf
en degrés $0$ et $-1$.
D'après Catégories dérivées d'espaces, lemme \ref{spaces-perfect-lemma-ext-finite},
les groupes $\Ext$ affichés sont des $k$-espaces vectoriels de dimension finie,
ce qui achève la démonstration.
\end{proof}

\noindent
Dans l'exemple \ref{example-spaces}, si $X$ est un espace algébrique propre sur $k$,
alors $\Deformationcategory_X$
admet une présentation par un groupoïde lisse proreprésentable en foncteurs
sur $\mathcal{C}_\Lambda$
et possède a fortiori un objet formel versel (minimal). Cela résulte
des lemmes \ref{lemma-spaces-RS} et
\ref{lemma-proper-spaces-TI},
ainsi que de la discussion générale de la section \ref{section-general}.

\begin{lemma}
\label{lemma-spaces-hull}
Dans l'exemple \ref{example-spaces}, supposons que $X$ soit un espace algébrique propre sur $k$.
Supposons que $\Lambda$ soit un anneau local complet de corps résiduel $k$
(cas classique). Alors le foncteur
$$
F : \mathcal{C}_\Lambda \longrightarrow \textit{Ensembles},\quad
A \longmapsto \Ob(\Deformationcategory_X(A))/\cong
$$
des classes d'isomorphisme d'objets possède une enveloppe. Si
$\text{Der}_k(\mathcal{O}_X, \mathcal{O}_X) = 0$, alors
$F$ est proreprésentable.
\end{lemma}

\begin{proof}
L'existence d'une enveloppe résulte immédiatement des
lemmes \ref{lemma-spaces-RS} et \ref{lemma-proper-spaces-TI}, ainsi que de
Théorie formelle des déformations, lemme \ref{formal-defos-lemma-RS-implies-S1-S2}
et remarque \ref{formal-defos-remark-compose-minimal-into-iso-classes}.

\medskip\noindent
Supposons $\text{Der}_k(\mathcal{O}_X, \mathcal{O}_X) = 0$. Alors
$\Deformationcategory_X$ et $F$ sont équivalents d'après
Théorie formelle des déformations, lemme \ref{formal-defos-lemma-infdef-trivial}.
Ainsi, $F$ est un foncteur de déformations (car $\Deformationcategory_X$ est une
catégorie de déformations) d'espace tangent de dimension finie, et nous pouvons appliquer
Théorie formelle des déformations, théorème
\ref{formal-defos-theorem-Schlessinger-prorepresentability}.
\end{proof}





\section{Déformations des complétés}
\label{section-compare}

\noindent
Dans cette section, nous comparons les problèmes de déformation posés
par une algèbre et par son complété.
Nous discutons d'abord de la « possibilité de relèvement ».

\begin{lemma}
\label{lemma-lift-equivalence-module-derived}
Soit $A' \to A$ une surjection d'anneaux de noyau nilpotent.
Soit $A' \to P'$ un homomorphisme d'anneaux plat.
Posons $P = P' \otimes_{A'} A$.
Soit $M$ un module plat sur $A$, muni d'une structure de $P$-module.
Alors les assertions suivantes sont équivalentes :
\begin{enumerate}
\item il existe un module plat sur $A'$, muni d'une structure de $P'$-module, noté $M'$, tel que
$M' \otimes_{P'} P = M$ ;
\item il existe un objet $K' \in D^-(P')$ tel que
$K' \otimes_{P'}^\mathbf{L} P = M$.
\end{enumerate}
\end{lemma}

\begin{proof}
Supposons que $M'$ soit comme dans (1). Alors
$$
M = M' \otimes_P P' = M' \otimes_{A'} A =
M' \otimes_A^\mathbf{L} A' = M' \otimes_{P'}^\mathbf{L} P
$$
Les deux premières égalités sont claires, la troisième vaut parce que
$M'$ est plat sur $A'$, et la quatrième résulte de
Compléments sur l'algèbre, lemme \ref{more-algebra-lemma-base-change-comparison}.
Ainsi, (2) est vérifiée. Réciproquement, supposons $K'$ comme dans (2).
Nous pouvons supposer, et supposons, $M$ non nul.
Soit $t$ le plus grand entier tel que $H^t(K')$ soit non nul
(il existe puisque $M$ est non nul).
Alors $H^t(K') \otimes_{P'} P = H^t(K' \otimes_{P'}^\mathbf{L} P)$
est nul si $t > 0$. Comme le noyau de $P' \to P$ est nilpotent,
cela implique $H^t(K') = 0$ par le lemme de Nakayama, contradiction.
Ainsi, $t = 0$ (le cas $t < 0$ est lui aussi absurde).
Alors $M' = H^0(K')$ est un $P'$-module tel que $M = M' \otimes_{P'} P$,
et la suite spectrale de Tor donne une application injective
$$
\text{Tor}_1^{P'}(M', P) \to H^{-1}(M' \otimes_{P'}^\mathbf{L} P) = 0
$$
Par la référence ci-dessus sur le changement de base dérivé,
$0 = \text{Tor}_1^{P'}(M', P) = \text{Tor}_1^{A'}(M', A)$.
Nous concluons que $M'$ est plat sur $A'$ d'après
Algèbre, lemme \ref{algebra-lemma-what-does-it-mean}.
\end{proof}

\begin{lemma}
\label{lemma-lift-equivalence-module}
Considérons un diagramme commutatif d'anneaux noethériens
$$
\xymatrix{
A' \ar[d] \ar[r] &
P' \ar[d] \ar[r] &
Q' \ar[d] \\
A \ar[r] &
P \ar[r] &
Q
}
$$
dont les carrés sont cartésiens, les flèches horizontales plates et les
flèches verticales surjectives, de noyaux nilpotents.
Soit $J' \subset P'$ un idéal tel que $P'/J' = Q'/J'Q'$.
Soit $M$ un module plat sur $A$, muni d'une structure de $P$-module.
Supposons que, pour tout $g \in J'$, il existe un module plat sur $A'$ muni d'une structure de $(P')_g$-module
relevant $M_g$. Alors les assertions suivantes sont équivalentes :
\begin{enumerate}
\item $M$ possède un relèvement plat sur $A'$ muni d'une structure de $P'$-module ;
\item $M \otimes_P Q$ possède un relèvement plat sur $A'$ muni d'une structure de $Q'$-module.
\end{enumerate}
\end{lemma}

\begin{proof}
Soit $I = \Ker(A' \to A)$. Par récurrence sur l'entier $n > 1$
tel que $I^n = 0$, nous nous ramenons au cas où $I$ est un idéal
de carré nul ; nous omettons les détails.
Nous traduisons la condition de possibilité de relèvement de
$M$ en le problème de trouver un objet de $D^-(P')$ comme dans le
lemme \ref{lemma-lift-equivalence-module-derived}.
L'obstruction à cette construction est l'élément
$$
\omega(M) \in \text{Ext}^2_P(M, M \otimes_P^\mathbf{L} IP) =
\text{Ext}^2_P(M, M \otimes_P IP)
$$
construit dans
Théorie des déformations, lemme \ref{defos-lemma-canonical-class-algebra}.
L'égalité dans la formule affichée vaut car
$M \otimes_P^\mathbf{L} IP = M \otimes_P IP$
puisque $M$ et $P$ sont plats sur $A$\footnote{Choisissons une résolution
$F_\bullet \to I$ par des $A$-modules libres. Puisque $A \to P$ est plat,
$P \otimes_A F_\bullet$ est une résolution libre de $IP$.
Ainsi, $M \otimes_P^\mathbf{L} IP$ est représenté par
$M \otimes_P P \otimes_A F_\bullet = M \otimes_A F_\bullet$.
Ce complexe n'a de cohomologie qu'en degré $0$, car $M$ est plat sur $A$.}.
De même, l'obstruction au relèvement de $M \otimes_P Q$ est
l'élément
$$
\omega(M \otimes_P Q) \in
\text{Ext}^2_Q(M \otimes_P Q, (M \otimes_P Q) \otimes_Q IQ)
$$
qui est l'image de $\omega(M)$ par la fonctorialité
de la construction $\omega(-)$ de
Théorie des déformations, lemme \ref{defos-lemma-canonical-class-algebra}.
D'après Compléments sur l'algèbre, lemme \ref{more-algebra-lemma-base-change-RHom},
nous avons
$$
\text{Ext}^2_Q(M \otimes_P Q, (M \otimes_P Q) \otimes_Q IQ) =
\text{Ext}^2_P(M, M \otimes_P IP) \otimes_P Q
$$
où nous utilisons le fait que $P$ est noethérien et $M$ fini.
Notre hypothèse sur $P' \to Q'$ garantit que, pour un $P$-module $E$,
l'application $E \to E \otimes_P Q$ est bijective sur la partie annulée par une puissance de $J'$ ; voir
Compléments sur l'algèbre, lemme
\ref{more-algebra-lemma-neighbourhood-equivalence}.
Nous concluons donc qu'il suffit de montrer que $\omega(M)$
est annulé par une puissance de $J'$. Autrement dit, il suffit de montrer que
$\omega(M)$ s'annule dans
$$
\text{Ext}^2_P(M, M \otimes_P IP)_g =
\text{Ext}^2_{P_g}(M_g, M_g \otimes_{P_g} IP_g)
$$
pour tout $g \in J'$. Cependant, par la compatibilité de la formation de $\omega(M)$
au changement de base, nous concluons à nouveau que cela est vrai puisque $M_g$
est supposé posséder un relèvement (il faut bien sûr utiliser de nouveau toute la
chaîne d'équivalences).
\end{proof}

\begin{lemma}
\label{lemma-lift-equivalence}
Soit $A' \to A$ un homomorphisme surjectif d'anneaux noethériens de noyau nilpotent.
Soit $A \to B$ un homomorphisme d'anneaux plat de type fini.
Soit $\mathfrak b \subset B$ un idéal tel que
$\Spec(B) \to \Spec(A)$ soit syntomique sur le complémentaire de $V(\mathfrak b)$.
Alors $B$ possède un relèvement plat sur $A'$ si et seulement si le complété
$\mathfrak b$-adique $B^\wedge$ possède un relèvement plat sur $A'$.
\end{lemma}

\begin{proof}
Choisissons une surjection de $A$-algèbres $P = A[x_1, \ldots, x_n] \to B$.
Soit $\mathfrak p \subset P$ l'image réciproque de $\mathfrak b$.
Posons $P' = A'[x_1, \ldots, x_n]$ et notons $\mathfrak p' \subset P'$
l'image réciproque de $\mathfrak p$. (Bien sûr, $\mathfrak p$
et $\mathfrak p'$ ne désignent pas ici des idéaux premiers.)
Nous noterons $P^\wedge$ et $(P')^\wedge$ les complétés respectifs.

\medskip\noindent
Supposons que $A' \to B'$ soit un relèvement plat de $A \to B$, autrement dit
que $A' \to B'$ soit plat et qu'il existe un isomorphisme de $A$-algèbres
$B = B' \otimes_{A'} A$. Nous pouvons alors choisir un homomorphisme de $A'$-algèbres
$P' \to B'$ relevant la surjection donnée $P \to B$.
Par le lemme de Nakayama (Algèbre, lemme \ref{algebra-lemma-NAK}),
nous trouvons que $B'$ est un quotient de $P'$. En particulier, nous trouvons
que nous pouvons rendre $B'$ plat sur $A'$ comme $P'$-module,
relevant $B$, qui est plat sur $A$, comme $P$-module.
Réciproquement, si nous pouvons relever $B$ en un $P'$-module $M'$ plat sur $A'$,
alors $M'$ est un module cyclique $M' \cong P'/J'$ (en utilisant de nouveau Nakayama),
et, en posant $B' = P'/J'$, nous obtenons un relèvement plat de $B$ comme algèbre.

\medskip\noindent
Posons $C = B^\wedge$ et $\mathfrak c = \mathfrak bC$.
Supposons que $A' \to C'$ soit un relèvement plat de $A \to C$.
Alors $C'$ est complet pour l'image réciproque $\mathfrak c'$
de $\mathfrak c$
(Algèbre, lemme \ref{algebra-lemma-complete-modulo-nilpotent}).
Choisissons un homomorphisme de $A'$-algèbres $P' \to C'$ relevant
l'homomorphisme de $A$-algèbres $P \to C$. Ces homomorphismes se prolongent aux
complétés et donnent des surjections $P^\wedge \to C$ et $(P')^\wedge \to C'$
(pour la seconde, en utilisant de nouveau le lemme de Nakayama).
En particulier, nous trouvons que $C'$ est plat sur $A'$ et peut être muni d'une structure de
$(P')^\wedge$-module relevant $C$, plat sur $A$, comme $P^\wedge$-module.
Réciproquement, si nous pouvons relever $C$ en un $(P')^\wedge$-module $N'$ plat sur $A'$,
alors $N'$ est un module cyclique $N' \cong (P')^\wedge/\tilde J$
(en utilisant de nouveau Nakayama) et, en posant $C' = (P')^\wedge/\tilde J$,
nous obtenons un relèvement plat de $C$ comme algèbre.

\medskip\noindent
Observons que $P' \to (P')^\wedge$ est un homomorphisme d'anneaux plat qui
induit un isomorphisme $P'/\mathfrak p' = (P')^\wedge/\mathfrak p'(P')^\wedge$.
Nous concluons que notre lemme résulte du
lemme \ref{lemma-lift-equivalence-module}, pourvu que nous puissions
montrer que $B_g$ se relève en un module plat sur $A'$ muni d'une structure de $P'_g$-module pour
$g \in \mathfrak p'$. Or l'homomorphisme d'anneaux $A \to B_g$ est syntomique
et se relève donc en une $A'$-algèbre plate $B'$ d'après
Lissage des homomorphismes d'anneaux, proposition \ref{smoothing-proposition-lift-smooth}.
Puisque $A' \to P'_g$ est lisse, nous pouvons relever $P_g \to B_g$
en une application surjective $P'_g \to B'$ comme ci-dessus, ce qui donne le résultat voulu.
\end{proof}

\noindent
Notation. Soit $A \to B$ un homomorphisme d'anneaux. Soit $N$ un $B$-module.
Nous notons $\text{Exal}_A(B, N)$ l'ensemble des classes d'isomorphisme
des extensions
$$
0 \to N \to C \to B \to 0
$$
de $A$-algèbres telles que $N$ soit un idéal de carré nul dans $C$.
Étant donnée une seconde extension $0 \to N \to C' \to B \to 0$ de ce type, un isomorphisme
est un isomorphisme de $A$-algèbres $C \to C'$ tel que le diagramme
$$
\xymatrix{
0 \ar[r] &
N \ar[r] \ar[d]_{\text{id}} &
C \ar[r] \ar[d] &
B \ar[r] \ar[d]_{\text{id}} & 0 \\
0 \ar[r] &
N \ar[r] &
C' \ar[r] &
B \ar[r] & 0
}
$$
soit commutatif. La correspondance $N \mapsto \text{Exal}_A(B, N)$
est un foncteur qui transforme les produits en produits.
C'est donc un foncteur additif, et $\text{Exal}_A(B, N)$
possède une structure naturelle de $B$-module. En fait, d'après
Théorie des déformations, lemme \ref{defos-lemma-choices},
nous avons $\text{Exal}_A(B, N) = \text{Ext}^1_B(\NL_{B/A}, N)$.

\begin{lemma}
\label{lemma-first-order-completion}
Soit $k$ un corps. Soit $B$ une $k$-algèbre de type fini.
Soit $J \subset B$ un idéal tel que
$\Spec(B) \to \Spec(k)$ soit lisse sur le complémentaire de $V(J)$.
Soit $N$ un $B$-module fini.
Alors il existe une bijection canonique
$$
\text{Exal}_k(B, N) \to \text{Exal}_k(B^\wedge, N^\wedge)
$$
Ici, $B^\wedge$ et $N^\wedge$ sont les complétés $J$-adiques.
\end{lemma}

\begin{proof}
L'application est donnée par complétion : étant donné $0 \to N \to C \to B \to 0$
dans $\text{Exal}_k(B, N)$, nous l'envoyons sur le complété $C^\wedge$
de $C$ pour l'image réciproque de $J$. Comparer avec
la démonstration du lemme \ref{lemma-completion}.

\medskip\noindent
Puisque $k \to B$ est de présentation finie, le complexe
$\NL_{B/k}$ peut être représenté par un complexe
$N^{-1} \to N^0$, où $N^i$ est un $B$-module fini ; voir
Algèbre, section \ref{algebra-section-netherlander}, et
en particulier
Algèbre, lemme \ref{algebra-lemma-NL-homotopy}.
Comme $B$ est noethérien, cela signifie que $\NL_{B/k}$
est pseudo-cohérent. Pour $g \in J$, la $k$-algèbre $B_g$
est lisse et donc $(\NL_{B/k})_g = \NL_{B_g/k}$
est quasi-isomorphe à un $B$-module projectif de type fini placé en degré $0$.
Ainsi, $\text{Ext}^i_B(\NL_{B/k}, N)_g = 0$ pour $i \geq 1$
et tout $B$-module $N$. D'après
Compléments sur l'algèbre, lemme \ref{more-algebra-lemma-ext-annihilated-into},
nous concluons que
$$
\text{Ext}^1_B(\NL_{B/k}, N) \longrightarrow
\lim_n \text{Ext}^1_B(\NL_{B/k}, N/J^n N)
$$
est un isomorphisme pour tout $B$-module fini $N$.

\medskip\noindent
Injectivité de l'application.
Supposons que $0 \to N \to C \to B \to 0$ appartienne à $\text{Exal}_k(B, N)$
et s'envoie sur zéro dans $\text{Exal}_k(B^\wedge, N^\wedge)$.
Choisissons un scindage $C^\wedge = B^\wedge \oplus N^\wedge$.
Alors l'application induite $C \to C^\wedge \to N^\wedge$
donne des applications $C \to N/J^nN$ pour tout $n$.
Nous voyons donc que notre élément appartient au noyau des applications
$$
\text{Ext}^1_B(\NL_{B/k}, N) \to
\text{Ext}^1_B(\NL_{B/k}, N/J^n N)
$$
pour tout $n$. Le paragraphe précédent montre que
notre élément est nul.

\medskip\noindent
Surjectivité de l'application. Soit $0 \to N^\wedge \to C' \to B^\wedge \to 0$
un élément de $\text{Exal}_k(B^\wedge, N^\wedge)$.
Par image réciproque selon $B \to B^\wedge$, nous obtenons un élément
$0 \to N^\wedge \to C'' \to B \to 0$ dans
$\text{Exal}_k(B, N^\wedge)$.
nous avons
$$
\text{Ext}^1_B(\NL_{B/k}, N^\wedge) =
\text{Ext}^1_B(\NL_{B/k}, N) \otimes_B B^\wedge =
\text{Ext}^1_B(\NL_{B/k}, N)
$$
La première égalité vaut puisque $N^\wedge = N \otimes_B B^\wedge$
(Algèbre, lemme \ref{algebra-lemma-completion-tensor})
et d'après
Compléments sur l'algèbre, lemme \ref{more-algebra-lemma-pseudo-coherence-and-ext}.
La seconde égalité vaut parce que $\text{Ext}^1_B(\NL_{B/k}, N)$
est annulé par une puissance de $J$ (voir ci-dessus), que $B \to B^\wedge$ est plat et induit
un isomorphisme $B/J \to B^\wedge/JB^\wedge$, et d'après
Compléments sur l'algèbre, lemme \ref{more-algebra-lemma-neighbourhood-equivalence}.
Nous pouvons donc trouver un $C \in \text{Exal}_k(B, N)$ s'envoyant sur $C''$ dans
$\text{Exal}_k(B, N^\wedge)$.
Ainsi,
$$
0 \to N^\wedge \to C' \to B^\wedge \to 0
\quad\text{et}\quad
0 \to N^\wedge \to C^\wedge \to B^\wedge \to 0
$$
ces suites sont deux éléments de $\text{Exal}_k(B^\wedge, N^\wedge)$
qui s'envoient sur le même élément de $\text{Exal}_k(B, N^\wedge)$.
En prenant leur différence, nous obtenons un élément
$0 \to N^\wedge \to C' \to B^\wedge \to 0$ de
$\text{Exal}_k(B^\wedge, N^\wedge)$
dont l'image dans $\text{Exal}_k(B, N^\wedge)$ est nulle.
Cela signifie qu'il existe
$$
\xymatrix{
0 \ar[r] &
N^\wedge \ar[r] &
C' \ar[r] &
B^\wedge \ar[r] & 0 \\
& & B \ar[u]^\sigma \ar[ru]
}
$$
Soit $J' \subset C'$ l'image réciproque de $JB^\wedge \subset B^\wedge$.
Pour achever la démonstration, il suffit de noter que
$\sigma$ est continue pour la topologie $J$-adique sur $B$
et la topologie $J'$-adique sur $C'$, et que $C'$ est complet pour la topologie $J'$-adique d'après
Algèbre, lemme \ref{algebra-lemma-complete-modulo-nilpotent}
(nous utilisons aussi ici que $C'$ est noethérien ; un petit détail est omis).
En effet, cela signifie que $\sigma$ se factorise par le
complété $B^\wedge$ et que $C' = 0$ dans $\text{Exal}_k(B^\wedge, N^\wedge)$.
\end{proof}

\begin{lemma}
\label{lemma-smooth-completion}
Dans l'exemple \ref{example-rings}, soit $P$ une $k$-algèbre.
Soit $J \subset P$ un idéal.
Notons $P^\wedge$ le complété $J$-adique. Si
\begin{enumerate}
\item $k \to P$ est de type fini ;
\item $\Spec(P) \to \Spec(k)$ est lisse sur le complémentaire de $V(J)$.
\end{enumerate}
alors le foncteur entre catégories de déformations du
lemme \ref{lemma-completion}
$$
\Deformationcategory_P \longrightarrow \Deformationcategory_{P^\wedge}
$$
est lisse et induit un isomorphisme sur les espaces tangents.
\end{lemma}

\begin{proof}
Nous savons que $\Deformationcategory_P$ et $\Deformationcategory_{P^\wedge}$
sont des catégories de déformations d'après le lemme \ref{lemma-rings-RS}.
Il suffit donc de vérifier que
notre foncteur identifie les espaces tangents et les possibilités
de relèvement ; voir Théorie formelle des déformations, lemme
\ref{formal-defos-lemma-easy-check-smooth}.
La propriété concernant les possibilités de relèvement est démontrée dans le
lemme \ref{lemma-lift-equivalence},
et l'isomorphisme sur les espaces tangents est le cas particulier du
lemme \ref{lemma-first-order-completion} où $N = B$.
\end{proof}



\section{Déformations des localisations}
\label{section-compare-localization}

\noindent
Dans cette section, nous comparons les problèmes de déformation posés
par une algèbre et par sa localisation en une partie multiplicative.
Nous discutons d'abord de la « possibilité de relèvement ».

\begin{lemma}
\label{lemma-lift-equivalence-localization}
Soit $A' \to A$ un homomorphisme surjectif d'anneaux noethériens de noyau nilpotent.
Soit $A \to B$ un homomorphisme d'anneaux plat de type fini.
Soit $S \subset B$ une partie multiplicative telle que,
si $\Spec(B) \to \Spec(A)$ n'est pas syntomique en $\mathfrak q$,
alors $S \cap \mathfrak q = \emptyset$.
Alors $B$ possède un relèvement plat sur $A'$ si et seulement si
$S^{-1}B$ possède un relèvement plat sur $A'$.
\end{lemma}

\begin{proof}
Cette démonstration est la même que celle du
lemme \ref{lemma-lift-equivalence}, mais plus facile. Nous suggérons au
lecteur de l'omettre.
Choisissons une surjection de $A$-algèbres $P = A[x_1, \ldots, x_n] \to B$.
Soit $S_P \subset P$ l'image réciproque de $S$.
Posons $P' = A'[x_1, \ldots, x_n]$ et notons $S_{P'} \subset P'$
l'image réciproque de $S_P$.

\medskip\noindent
Supposons que $A' \to B'$ soit un relèvement plat de $A \to B$, autrement dit
que $A' \to B'$ soit plat et qu'il existe un isomorphisme de $A$-algèbres
$B = B' \otimes_{A'} A$. Nous pouvons alors choisir un homomorphisme de $A'$-algèbres
$P' \to B'$ relevant la surjection donnée $P \to B$.
Par le lemme de Nakayama (Algèbre, lemme \ref{algebra-lemma-NAK}),
nous trouvons que $B'$ est un quotient de $P'$. En particulier, nous trouvons
que nous pouvons rendre $B'$ plat sur $A'$ comme $P'$-module,
relevant $B$, qui est plat sur $A$, comme $P$-module.
Réciproquement, si nous pouvons relever $B$ en un $P'$-module $M'$ plat sur $A'$,
alors $M'$ est un module cyclique $M' \cong P'/J'$ (en utilisant de nouveau Nakayama),
et, en posant $B' = P'/J'$, nous obtenons un relèvement plat de $B$ comme algèbre.

\medskip\noindent
Posons $C = S^{-1}B$. Supposons que $A' \to C'$ soit un relèvement plat de $A \to C$.
Les éléments de $C'$ dont l'image dans $C$ est inversible sont inversibles.
Choisissons un homomorphisme de $A'$-algèbres $P' \to C'$ relevant
l'homomorphisme de $A$-algèbres $P \to C$. D'après la remarque précédente,
ces homomorphismes se prolongent aux localisations et donnent des surjections
$S_P^{-1}P \to C$ et $S_{P'}^{-1}P' \to C'$
(pour la seconde, utiliser le lemme de Nakayama).
En particulier, nous trouvons que $C'$ est plat sur $A'$ et peut être muni d'une structure de
$S_{P'}^{-1}P'$-module relevant $C$, plat sur $A$, comme
$S_P^{-1}P$-module. Réciproquement, si nous pouvons relever $C$ en un
$S_{P'}^{-1}P'$-module $N'$ plat sur $A'$, alors $N'$
est un module cyclique $N' \cong S_{P'}^{-1}P'/\tilde J$
(en utilisant de nouveau Nakayama) et, en posant $C' = S_{P'}^{-1}P'/\tilde J$,
nous obtenons un relèvement plat de $C$ comme algèbre.

\medskip\noindent
Le lieu syntomique d'un morphisme de schémas est ouvert par définition.
Soit $J_B \subset B$ un idéal définissant l'ensemble des points
de $\Spec(B)$ où $\Spec(B) \to \Spec(A)$ n'est pas syntomique.
Notons $J_P \subset P$ et $J_{P'} \subset P'$ les idéaux
correspondants. Observons que $P' \to S_{P'}^{-1}P'$ est un homomorphisme d'anneaux plat qui
induit un isomorphisme $P'/J_{P'} = S_{P'}^{-1}P'/J_{P'}S_{P'}^{-1}P'$
grâce à notre hypothèse sur $S$ dans le lemme ; plus précisément, l'hypothèse
du lemme signifie exactement que $B/J_B = S^{-1}(B/J_B)$.
Nous concluons que notre lemme résulte du
lemme \ref{lemma-lift-equivalence-module}, pourvu que nous puissions
montrer que $B_g$ se relève en un module plat sur $A'$ muni d'une structure de $P'_g$-module pour
$g \in J_B$. Or l'homomorphisme d'anneaux $A \to B_g$ est syntomique
et se relève donc en une $A'$-algèbre plate $B'$ d'après
Lissage des homomorphismes d'anneaux, proposition \ref{smoothing-proposition-lift-smooth}.
Puisque $A' \to P'_g$ est lisse, nous pouvons relever $P_g \to B_g$
en une application surjective $P'_g \to B'$ comme ci-dessus, ce qui donne le résultat voulu.
\end{proof}

\begin{lemma}
\label{lemma-first-order-localization}
Soit $k$ un corps. Soit $B$ une $k$-algèbre de type fini.
Soit $S \subset B$ une partie multiplicative telle que,
si $\Spec(B) \to \Spec(k)$ n'est pas lisse en $\mathfrak q$,
alors $S \cap \mathfrak q = \emptyset$.
Soit $N$ un $B$-module fini.
Alors il existe une bijection canonique
$$
\text{Exal}_k(B, N) \to \text{Exal}_k(S^{-1}B, S^{-1}N)
$$
\end{lemma}

\begin{proof}
Cette démonstration est la même que celle du
lemme \ref{lemma-first-order-completion}, mais plus facile. Nous suggérons au
lecteur de l'omettre.
L'application est donnée par localisation : étant donné $0 \to N \to C \to B \to 0$
dans $\text{Exal}_k(B, N)$, nous l'envoyons sur la localisation $S_C^{-1}C$
de $C$ en l'image réciproque $S_C \subset C$ de $S$.
Comparer avec la démonstration du lemme \ref{lemma-localization}.

\medskip\noindent
Le lieu de lissité d'un morphisme de schémas est ouvert par définition.
Soit $J \subset B$ un idéal définissant l'ensemble des points
de $\Spec(B)$ où $\Spec(B) \to \Spec(A)$ n'est pas lisse.
Puisque $k \to B$ est de présentation finie, le complexe
$\NL_{B/k}$ peut être représenté par un complexe
$N^{-1} \to N^0$, où $N^i$ est un $B$-module fini ; voir
Algèbre, section \ref{algebra-section-netherlander}, et
en particulier
Algèbre, lemme \ref{algebra-lemma-NL-homotopy}.
Comme $B$ est noethérien, cela signifie que $\NL_{B/k}$
est pseudo-cohérent. Pour $g \in J$, la $k$-algèbre $B_g$
est lisse et donc $(\NL_{B/k})_g = \NL_{B_g/k}$
est quasi-isomorphe à un $B$-module projectif de type fini placé en degré $0$.
Ainsi, $\text{Ext}^i_B(\NL_{B/k}, N)_g = 0$ pour $i \geq 1$
et tout $B$-module $N$. Enfin, nous avons
$$
\text{Ext}^1_{S^{-1}B}(\NL_{S^{-1}B/k}, S^{-1}N) =
\text{Ext}^1_B(\NL_{B/k}, N) \otimes_B S^{-1}B =
\text{Ext}^1_B(\NL_{B/k}, N)
$$
La première égalité résulte de
Compléments sur l'algèbre, lemme \ref{more-algebra-lemma-base-change-RHom},
et d'Algèbre, lemme \ref{algebra-lemma-localize-NL}.
La seconde vaut parce que $\text{Ext}^1_B(\NL_{B/k}, N)$ est annulé par une puissance de $J$
et que les éléments de $S$ agissent inversiblement sur les modules annulés par une puissance de $J$.
Cela achève la démonstration grâce à la description de $\text{Exal}_A(B, N)$
comme $\text{Ext}^1_B(\NL_{B/A}, N)$ donnée juste avant le
lemme \ref{lemma-first-order-completion}.
\end{proof}

\begin{lemma}
\label{lemma-smooth-localization}
Dans l'exemple \ref{example-rings}, soit $P$ une $k$-algèbre.
Soit $S \subset P$ une partie multiplicative. Si
\begin{enumerate}
\item $k \to P$ est de type fini ;
\item $\Spec(P) \to \Spec(k)$ est lisse en tout point de
$V(g)$ pour tout $g \in S$.
\end{enumerate}
alors le foncteur entre catégories de déformations du
lemme \ref{lemma-localization}
$$
\Deformationcategory_P \longrightarrow \Deformationcategory_{S^{-1}P}
$$
est lisse et induit un isomorphisme sur les espaces tangents.
\end{lemma}

\begin{proof}
Nous savons que $\Deformationcategory_P$ et $\Deformationcategory_{S^{-1}P}$
sont des catégories de déformations d'après le lemme \ref{lemma-rings-RS}.
Il suffit donc de vérifier que
notre foncteur identifie les espaces tangents et les possibilités
de relèvement ; voir Théorie formelle des déformations, lemme
\ref{formal-defos-lemma-easy-check-smooth}.
La propriété concernant les possibilités de relèvement est démontrée dans le
lemme \ref{lemma-lift-equivalence-localization},
et l'isomorphisme sur les espaces tangents est le cas particulier du
lemme \ref{lemma-first-order-localization} où $N = B$.
\end{proof}



\section{Déformations des hensélisés}
\label{section-compare-henselization}

\noindent
Dans cette section, nous comparons les problèmes de déformation posés
par une algèbre et par son complété.
Nous discutons d'abord de la « possibilité de relèvement ».

\begin{lemma}
\label{lemma-lift-equivalence-henselization}
Soit $A' \to A$ un homomorphisme surjectif d'anneaux noethériens de noyau nilpotent.
Soit $A \to B$ un homomorphisme d'anneaux plat de type fini.
Soit $\mathfrak b \subset B$ un idéal tel que
$\Spec(B) \to \Spec(A)$ soit syntomique sur le complémentaire de $V(\mathfrak b)$.
Soit $(B^h, \mathfrak b^h)$ l'hensélisation du couple $(B, \mathfrak b)$.
Alors $B$ possède un relèvement plat sur $A'$ si et seulement si $B^h$ possède un relèvement plat sur $A'$.
\end{lemma}

\begin{proof}[Première démonstration]
Cette démonstration triche quelque peu. En effet, si $B$ possède un relèvement plat $B'$, alors,
en prenant l'hensélisé $(B')^h$, nous obtenons un relèvement plat de $B^h$
(comparer avec la démonstration du lemme \ref{lemma-henselization}).
Réciproquement, supposons que $C'$ soit un relèvement plat sur $A'$ de $(B')^h$.
Soit alors $\mathfrak c' \subset C'$ l'image réciproque de
l'idéal $\mathfrak b^h$. Le complété $(C')^\wedge$ de
$C'$ pour $\mathfrak c'$ est alors un relèvement de $B^\wedge$ (détails omis).
Nous voyons donc que $B$ possède un relèvement plat d'après le
lemme \ref{lemma-lift-equivalence}.
\end{proof}

\begin{proof}[Seconde démonstration]
Choisissons une surjection de $A$-algèbres $P = A[x_1, \ldots, x_n] \to B$.
Soit $\mathfrak p \subset P$ l'image réciproque de $\mathfrak b$.
Posons $P' = A'[x_1, \ldots, x_n]$ et notons $\mathfrak p' \subset P'$
l'image réciproque de $\mathfrak p$. (Bien sûr, $\mathfrak p$
et $\mathfrak p'$ ne désignent pas ici des idéaux premiers.)
Nous noterons $P^h$ et $(P')^h$ les hensélisés respectifs.
Nous utiliserons la fonctorialité de l'hensélisation ainsi que le fait que
l'hensélisé d'un quotient est le quotient correspondant
de l'hensélisé ; voir
Compléments sur l'algèbre, lemmes
\ref{more-algebra-lemma-irreducible-henselian-pair-connected} et
\ref{more-algebra-lemma-henselization-integral}.

\medskip\noindent
Supposons que $A' \to B'$ soit un relèvement plat de $A \to B$, autrement dit
que $A' \to B'$ soit plat et qu'il existe un isomorphisme de $A$-algèbres
$B = B' \otimes_{A'} A$. Nous pouvons alors choisir un homomorphisme de $A'$-algèbres
$P' \to B'$ relevant la surjection donnée $P \to B$.
Par le lemme de Nakayama (Algèbre, lemme \ref{algebra-lemma-NAK}),
nous trouvons que $B'$ est un quotient de $P'$. En particulier, nous trouvons
que nous pouvons rendre $B'$ plat sur $A'$ comme $P'$-module,
relevant $B$, qui est plat sur $A$, comme $P$-module.
Réciproquement, si nous pouvons relever $B$ en un $P'$-module $M'$ plat sur $A'$,
alors $M'$ est un module cyclique $M' \cong P'/J'$ (en utilisant de nouveau Nakayama),
et, en posant $B' = P'/J'$, nous obtenons un relèvement plat de $B$ comme algèbre.

\medskip\noindent
Posons $C = B^h$ et $\mathfrak c = \mathfrak bC$.
Supposons que $A' \to C'$ soit un relèvement plat de $A \to C$.
Alors $C'$ est hensélien pour l'image réciproque
$\mathfrak c'$ de $\mathfrak c$
(d'après Compléments sur l'algèbre, lemme \ref{more-algebra-lemma-henselian-henselian-pair},
et parce que le noyau de $C' \to C$ est nilpotent).
Choisissons un homomorphisme de $A'$-algèbres $P' \to C'$ relevant
l'homomorphisme de $A$-algèbres $P \to C$. Ces homomorphismes se prolongent aux
hensélisés et donnent des surjections $P^h \to C$ et $(P')^h \to C'$
(pour la seconde, en utilisant de nouveau le lemme de Nakayama).
En particulier, nous trouvons que $C'$ est plat sur $A'$ et peut être muni d'une structure de
$(P')^h$-module relevant $C$, plat sur $A$, comme $P^h$-module.
Réciproquement, si nous pouvons relever $C$ en un $(P')^h$-module $N'$ plat sur $A'$,
alors $N'$ est un module cyclique $N' \cong (P')^h/\tilde J$
(en utilisant de nouveau Nakayama) et, en posant $C' = (P')^h/\tilde J$,
nous obtenons un relèvement plat de $C$ comme algèbre.

\medskip\noindent
Observons que $P' \to (P')^h$ est un homomorphisme d'anneaux plat qui
induit un isomorphisme $P'/\mathfrak p' = (P')^h/\mathfrak p'(P')^h$
(Compléments sur l'algèbre, lemme \ref{more-algebra-lemma-henselization-flat}).
Nous concluons que notre lemme résulte du
lemme \ref{lemma-lift-equivalence-module}, pourvu que nous puissions
montrer que $B_g$ se relève en un module plat sur $A'$ muni d'une structure de $P'_g$-module pour
$g \in \mathfrak p'$. Or l'homomorphisme d'anneaux $A \to B_g$ est syntomique
et se relève donc en une $A'$-algèbre plate $B'$ d'après
Lissage des homomorphismes d'anneaux, proposition \ref{smoothing-proposition-lift-smooth}.
Puisque $A' \to P'_g$ est lisse, nous pouvons relever $P_g \to B_g$
en une application surjective $P'_g \to B'$ comme ci-dessus, ce qui donne le résultat voulu.
\end{proof}

\begin{lemma}
\label{lemma-first-order-henselization}
Soit $k$ un corps. Soit $B$ une $k$-algèbre de type fini.
Soit $J \subset B$ un idéal tel que
$\Spec(B) \to \Spec(k)$ soit lisse sur le complémentaire de $V(J)$.
Soit $N$ un $B$-module fini.
Alors il existe une bijection canonique
$$
\text{Exal}_k(B, N) \to \text{Exal}_k(B^h, N^h)
$$
Ici, $(B^h, J^h)$ est l'hensélisation de $(B, J)$
et $N^h = N \otimes_B B^h$.
\end{lemma}

\begin{proof}
Cette démonstration est la même que celle du
lemme \ref{lemma-first-order-completion}, mais plus facile. Nous suggérons au
lecteur de l'omettre.
L'application est donnée par hensélisation : étant donné $0 \to N \to C \to B \to 0$
dans $\text{Exal}_k(B, N)$, nous l'envoyons sur
l'hensélisé $C^h$
de $C$ pour l'image réciproque $J_C \subset C$ de $J$.
Comparer avec la démonstration du lemme \ref{lemma-henselization}.

\medskip\noindent
Puisque $k \to B$ est de présentation finie, le complexe
$\NL_{B/k}$ peut être représenté par un complexe
$N^{-1} \to N^0$, où $N^i$ est un $B$-module fini ; voir
Algèbre, section \ref{algebra-section-netherlander}, et
en particulier
Algèbre, lemme \ref{algebra-lemma-NL-homotopy}.
Comme $B$ est noethérien, cela signifie que $\NL_{B/k}$
est pseudo-cohérent. Pour $g \in J$, la $k$-algèbre $B_g$
est lisse et donc $(\NL_{B/k})_g = \NL_{B_g/k}$
est quasi-isomorphe à un $B$-module projectif de type fini placé en degré $0$.
Ainsi, $\text{Ext}^i_B(\NL_{B/k}, N)_g = 0$ pour $i \geq 1$
et tout $B$-module $N$. Enfin, nous avons
\begin{align*}
\text{Ext}^1_{B^h}(\NL_{B^h/k}, N^h)
& =
\text{Ext}^1_{B^h}(\NL_{B/k} \otimes_B B^h, N \otimes_B B^h) \\
& =
\text{Ext}^1_B(\NL_{B/k}, N) \otimes_B B^h \\
& =
\text{Ext}^1_B(\NL_{B/k}, N)
\end{align*}
La première égalité résulte de
Compléments sur l'algèbre, lemme \ref{more-algebra-lemma-henselization-NL}
(ou plutôt de son analogue pour les hensélisations de couples).
La deuxième résulte de
Compléments sur l'algèbre, lemme \ref{more-algebra-lemma-base-change-RHom}.
La troisième vaut parce que $\text{Ext}^1_B(\NL_{B/k}, N)$ est annulé par une puissance de $J$,
que l'application $B \to B^h$ est plate et induit un isomorphisme
$B/J \to B^h/JB^h$ (Compléments sur l'algèbre, lemme
\ref{more-algebra-lemma-henselization-flat}), et d'après
Compléments sur l'algèbre, lemme \ref{more-algebra-lemma-neighbourhood-equivalence}.
Cela achève la démonstration grâce à la description de $\text{Exal}_A(B, N)$
comme $\text{Ext}^1_B(\NL_{B/A}, N)$ donnée juste avant le
lemme \ref{lemma-first-order-completion}.
\end{proof}

\begin{lemma}
\label{lemma-smooth-henselization}
Dans l'exemple \ref{example-rings}, soit $P$ une $k$-algèbre.
Soit $J \subset P$ un idéal.
Notons $(P^h, J^h)$ l'hensélisation du couple $(P, J)$. Si
\begin{enumerate}
\item $k \to P$ est de type fini ;
\item $\Spec(P) \to \Spec(k)$ est lisse sur le complémentaire de $V(J)$,
\end{enumerate}
alors le foncteur entre catégories de déformations du
lemme \ref{lemma-henselization}
$$
\Deformationcategory_P \longrightarrow \Deformationcategory_{P^h}
$$
est lisse et induit un isomorphisme sur les espaces tangents.
\end{lemma}

\begin{proof}
Nous savons que $\Deformationcategory_P$ et $\Deformationcategory_{P^h}$
sont des catégories de déformations d'après le lemme \ref{lemma-rings-RS}.
Il suffit donc de vérifier que
notre foncteur identifie les espaces tangents et les possibilités
de relèvement ; voir Théorie formelle des déformations, lemme
\ref{formal-defos-lemma-easy-check-smooth}.
La propriété concernant les possibilités de relèvement est démontrée dans le
lemme \ref{lemma-lift-equivalence-henselization},
et l'isomorphisme sur les espaces tangents est le cas particulier du
lemme \ref{lemma-first-order-henselization} où $N = B$.
\end{proof}



\section{Application aux singularités isolées}
\label{section-isolated}

\noindent
Nous appliquons la discussion précédente à l'étude de la théorie des déformations
d'une algèbre de type fini qui ne possède qu'un nombre fini de points singuliers.

\begin{lemma}
\label{lemma-isolated}
Dans l'exemple \ref{example-rings}, soit $P$ une $k$-algèbre.
Supposons que $k \to P$ soit de type fini et que $\Spec(P) \to \Spec(k)$
soit lisse sauf aux idéaux maximaux
$\mathfrak m_1, \ldots, \mathfrak m_n$ de $P$.
Soient $P_{\mathfrak m_i}$, $P_{\mathfrak m_i}^h$, $P_{\mathfrak m_i}^\wedge$
l'anneau local, l'hensélisé et le complété.
Alors les applications de catégories de déformations
$$
\Deformationcategory_P \to
\prod \Deformationcategory_{P_{\mathfrak m_i}} \to
\prod \Deformationcategory_{P_{\mathfrak m_i}^h} \to
\prod \Deformationcategory_{P_{\mathfrak m_i}^\wedge}
$$
sont lisses et induisent des isomorphismes sur leurs espaces tangents
de dimension finie.
\end{lemma}

\begin{proof}
L'espace tangent est de dimension finie d'après le
lemme \ref{lemma-finite-type-rings-TI}.
Les foncteurs entre les catégories sont construits
dans les lemmes \ref{lemma-localization}, \ref{lemma-henselization} et
\ref{lemma-completion} (nous omettons certaines vérifications du type :
le complété de l'hensélisé est le complété).

\medskip\noindent
Posons $J = \mathfrak m_1 \cap \ldots \cap \mathfrak m_n$ et appliquons le
lemme \ref{lemma-smooth-completion} pour obtenir que
$\Deformationcategory_P \to \Deformationcategory_{P^\wedge}$
est lisse et induit un isomorphisme sur les espaces tangents,
où $P^\wedge$ est le complété $J$-adique de $P$.
Cependant, puisque $P^\wedge = \prod P_{\mathfrak m_i}^\wedge$,
nous voyons que l'application $\Deformationcategory_P \to
\prod \Deformationcategory_{P_{\mathfrak m_i}^\wedge}$
est lisse et induit un isomorphisme sur les espaces tangents.

\medskip\noindent
Soit $(P^h, J^h)$ l'hensélisation du couple $(P, J)$.
Alors $P^h = \prod P_{\mathfrak m_i}^h$ (considérer les idempotents
et utiliser Compléments sur l'algèbre, lemme
\ref{more-algebra-lemma-characterize-henselian-pair}).
Nous pouvons donc appliquer le lemme \ref{lemma-smooth-henselization}
et conclure comme dans le cas du complété.

\medskip\noindent
Pour obtenir le dernier cas, il suffit de montrer que
$\Deformationcategory_{P_{\mathfrak m_i}} \to
\Deformationcategory_{P_{\mathfrak m_i}^\wedge}$
est lisse et induit des isomorphismes sur les espaces tangents pour chaque $i$ séparément.
Pour cela, nous pouvons remplacer $P$ par une localisation principale
dont l'unique point singulier est un idéal maximal $\mathfrak m$
(correspondant à $\mathfrak m_i$ dans le $P$ initial).
Nous pouvons alors appliquer le
lemme \ref{lemma-smooth-localization}
à la partie multiplicative $S = P \setminus \mathfrak m$ pour conclure.
Les détails mineurs sont omis.
\end{proof}






\section{Problèmes de déformation non obstrués}
\label{section-unobstructed}

\noindent
Soit $p : \mathcal{F} \to \mathcal{C}_\Lambda$ une
catégorie cofibrée en groupoïdes. Rappelons que nous disons que $\mathcal{F}$
est {\it lisse} ou {\it non obstruée} si $p$ est lisse.
Cela signifie qu'étant données une surjection $\varphi : A' \to A$ dans
$\mathcal{C}_\Lambda$ et $x \in \Ob(\mathcal{F}(A))$,
il existe un morphisme $f : x' \to x$ dans $\mathcal{F}$
tel que $p(f) = \varphi$.
Voir Théorie formelle des déformations, section \ref{formal-defos-section-smooth}.
Dans cette section, nous donnons quelques exemples géométriquement significatifs.

\begin{lemma}
\label{lemma-lci-unobstructed}
Dans l'exemple \ref{example-rings}, soit $P$ une intersection complète
locale sur $k$ (Algèbre, définition \ref{algebra-definition-lci-field}).
Alors $\Deformationcategory_P$ est non obstruée.
\end{lemma}

\begin{proof}
Soit $(A, Q) \to (k, P)$ un objet de $\Deformationcategory_P$.
Alors nous voyons que $A \to Q$ est un homomorphisme d'anneaux syntomique d'après
Algèbre, définition \ref{algebra-definition-lci}.
Ainsi, pour toute surjection $A' \to A$ dans $\mathcal{C}_\Lambda$,
nous voyons qu'il existe un morphisme $(A', Q') \to (A, Q)$
relevant $A' \to A$ d'après
Lissage des homomorphismes d'anneaux, proposition \ref{smoothing-proposition-lift-smooth}.
Cela démontre le lemme.
\end{proof}

\begin{lemma}
\label{lemma-glueing-smooth}
Dans la situation \ref{situation-glueing}, si $U_{12} \to \Spec(k)$ est lisse,
alors le morphisme
$$
\Deformationcategory_X
\longrightarrow
\Deformationcategory_{U_1} \times \Deformationcategory_{U_2} =
\Deformationcategory_{P_1} \times \Deformationcategory_{P_2}
$$
est lisse. Si, de plus,
$U_1$ est une intersection complète locale sur $k$, alors
$$
\Deformationcategory_X
\longrightarrow
\Deformationcategory_{U_2} = \Deformationcategory_{P_2}
$$
est lisse.
\end{lemma}

\begin{proof}
Les signes d'égalité sont justifiés par le lemme \ref{lemma-affine}.
Considérons $\mathcal{C}_\Lambda$ comme une catégorie de
déformations sur $\mathcal{C}_\Lambda$, comme dans
Théorie formelle des déformations, section \ref{formal-defos-section-smooth}.
Alors
$$
\Deformationcategory_{P_1} \times \Deformationcategory_{P_2} =
\Deformationcategory_{P_1}
\times_{\mathcal{C}_\Lambda}
\Deformationcategory_{P_2},
$$
voir Théorie formelle des déformations, remarques
\ref{formal-defos-remarks-cofibered-groupoids}
(\ref{formal-defos-item-product}).
En utilisant le
lemme \ref{lemma-glueing},
la première assertion revient à dire que le foncteur
$$
\Deformationcategory_{P_1}
\times_{\Deformationcategory_{P_{12}}}
\Deformationcategory_{P_2}
\longrightarrow
\Deformationcategory_{P_1}
\times_{\mathcal{C}_\Lambda}
\Deformationcategory_{P_2}
$$
est lisse. Cela résulte de Théorie formelle des déformations, lemme
\ref{formal-defos-lemma-map-fibre-products-smooth}, pourvu que
nous puissions montrer que $T\Deformationcategory_{P_{12}} = (0)$.
Cette annulation résulte du lemme \ref{lemma-smooth},
puisque $P_{12}$ est lisse sur $k$.
Pour la seconde assertion, il suffit de montrer que
$\Deformationcategory_{P_1} \to \mathcal{C}_\Lambda$
est lisse ; voir Théorie formelle des déformations, lemme
\ref{formal-defos-lemma-smooth-properties}.
Autrement dit, il faut montrer que $\Deformationcategory_{P_1}$
est non obstruée, ce qui est le lemme \ref{lemma-lci-unobstructed}.
\end{proof}

\begin{lemma}
\label{lemma-curve-isolated}
Dans l'exemple \ref{example-schemes}, soit $X$ un schéma sur $k$. Supposons que
\begin{enumerate}
\item $X$ soit séparé, de type fini sur $k$ et $\dim(X) \leq 1$ ;
\item $X \to \Spec(k)$ soit lisse sauf aux
points fermés $p_1, \ldots, p_n \in X$.
\end{enumerate}
Soient $\mathcal{O}_{X, p_1}$, $\mathcal{O}_{X, p_1}^h$,
$\mathcal{O}_{X, p_1}^\wedge$ l'anneau local, l'hensélisé et le complété.
Considérons les applications de catégories de déformations
$$
\Deformationcategory_X
\longrightarrow
\prod \Deformationcategory_{\mathcal{O}_{X, p_i}}
\longrightarrow
\prod \Deformationcategory_{\mathcal{O}_{X, p_i}^h}
\longrightarrow
\prod \Deformationcategory_{\mathcal{O}_{X, p_i}^\wedge}
$$
La première flèche est lisse, et les deuxième et troisième flèches
sont lisses et induisent des isomorphismes sur les espaces tangents.
\end{lemma}

\begin{proof}
Choisissons un ouvert affine $U_2 \subset X$ contenant
$p_1, \ldots, p_n$ et le point générique de toute composante
irréductible de $X$. Cela est possible d'après
Variétés, lemme \ref{varieties-lemma-dim-1-quasi-projective},
et Propriétés, lemme \ref{properties-lemma-ample-finite-set-in-affine}.
Alors $X \setminus U_2$ est fini, et nous pouvons choisir un ouvert affine
$U_1 \subset X \setminus \{p_1, \ldots, p_n\}$ tel que
$X = U_1 \cup U_2$. Posons $U_{12} = U_1 \cap U_2$.
Alors $U_1$ et $U_{12}$ sont des schémas affines lisses sur $k$.
Nous concluons que
$$
\Deformationcategory_X \longrightarrow \Deformationcategory_{U_2}
$$
est lisse d'après le lemme \ref{lemma-glueing-smooth}.
Le résultat découle des lemmes \ref{lemma-affine} et \ref{lemma-isolated}.
\end{proof}

\begin{lemma}
\label{lemma-curve-isolated-lci}
Dans l'exemple \ref{example-schemes}, soit $X$ un schéma sur $k$. Supposons que
\begin{enumerate}
\item $X$ soit séparé, de type fini sur $k$ et $\dim(X) \leq 1$ ;
\item $X$ soit une intersection complète locale sur $k$ ;
\item $X \to \Spec(k)$ soit lisse sauf en un nombre fini de points.
\end{enumerate}
Alors $\Deformationcategory_X$ est non obstruée.
\end{lemma}

\begin{proof}
Soient $p_1, \ldots, p_n \in X$ les points où $X \to \Spec(k)$
n'est pas lisse. Choisissons un ouvert affine $U_2 \subset X$ contenant
$p_1, \ldots, p_n$ et le point générique de toute composante
irréductible de $X$. Cela est possible d'après
Variétés, lemme \ref{varieties-lemma-dim-1-quasi-projective},
et Propriétés, lemme \ref{properties-lemma-ample-finite-set-in-affine}.
Alors $X \setminus U_2$ est fini, et nous pouvons choisir un ouvert affine
$U_1 \subset X \setminus \{p_1, \ldots, p_n\}$ tel que
$X = U_1 \cup U_2$. Posons $U_{12} = U_1 \cap U_2$.
Alors $U_1$ et $U_{12}$ sont des schémas affines lisses sur $k$.
Nous concluons que
$$
\Deformationcategory_X \longrightarrow \Deformationcategory_{U_2}
$$
est lisse d'après le lemme \ref{lemma-glueing-smooth}.
Le résultat découle des lemmes \ref{lemma-affine} et \ref{lemma-lci-unobstructed}.
\end{proof}




\section{Lissages}
\label{section-smoothing}

\noindent
Soit donné un schéma ou espace algébrique de type fini $X$ sur un corps $k$.
Il est souvent utile de trouver un morphisme plat de type fini $Y \to \Spec(k[[t]])$
dont la fibre générique est lisse et la fibre spéciale isomorphe à $X$.
Un tel objet est appelé un lissage de $X$. Dans cette section, nous trouverons
un lissage en dimension $1$ pour un espace séparé $X$ qui possède des singularités
isolées d'intersection complète locale.

\begin{lemma}
\label{lemma-criterion-smoothing}
Soit $k$ un corps. Posons $S = \Spec(k[[t]])$ et
$S_n = \Spec(k[t]/(t^n))$. Soit $Y \to S$ un morphisme propre et plat
de schémas dont la fibre spéciale $X$ est de Cohen-Macaulay et
équidimensionnelle de dimension $d$. Notons $X_n = Y \times_S S_n$.
Si, pour un certain $n \geq 1$, le $d$-ième idéal de Fitting de $\Omega_{X_n/S_n}$
contient $t^{n - 1}$, alors la fibre générique de $Y \to S$ est lisse.
\end{lemma}

\begin{proof}
D'après Compléments sur les morphismes, lemme
\ref{more-morphisms-lemma-flat-finite-presentation-CM-open}
nous voyons que $Y \to S$ est un morphisme de Cohen-Macaulay.
D'après Morphismes, lemme
\ref{morphisms-lemma-flat-finite-presentation-CM-fibres-relative-dimension}
nous voyons que $Y \to S$ est de dimension relative $d$.
D'après Diviseurs, lemme \ref{divisors-lemma-d-fitting-ideal-omega-smooth},
le $d$-ième idéal de Fitting $\mathcal{I} \subset \mathcal{O}_Y$
de $\Omega_{Y/S}$ définit le lieu singulier du morphisme $Y \to S$.
Autrement dit, $V(\mathcal{I}) \subset Y$ est le sous-ensemble fermé
des points où $Y \to S$ n'est pas lisse.
D'après Diviseurs, lemme \ref{divisors-lemma-base-change-and-fitting-ideal-omega},
la formation de cet idéal de Fitting commute au changement de base.
Par hypothèse, nous voyons que $t^{n - 1}$ est une section de
$\mathcal{I} + t^n\mathcal{O}_Y$. Ainsi, pour tout
$x \in X = V(t) \subset Y$, nous concluons que
$t^{n - 1} \in \mathcal{I}_x$, où $\mathcal{I}_x$ est le germe en $x$.
Cela implique que $V(\mathcal{I}) \subset V(t)$ dans un
voisinage ouvert de $X$ dans $Y$. Puisque $Y \to S$
est propre, il s'ensuit que $V(\mathcal{I}) \subset V(t)$,
comme souhaité.
\end{proof}

\begin{lemma}
\label{lemma-jouanolou-type-thing}
Soit $k$ un corps. Soient $1 \leq c \leq n$ des entiers.
Soient $f_1, \ldots, f_c \in k[x_1, \ldots x_n]$ des éléments.
Soient $a_{ij}$, $0 \leq i \leq n$, $1 \leq j \leq c$ des
indéterminées. Considérons
$$
g_j = f_j + a_{0j} + a_{1j}x_1 + \ldots + a_{nj}x_n \in
k[a_{ij}][x_1, \ldots, x_n]
$$
Notons $Y \subset \mathbf{A}^{n + c(n + 1)}_k$
le sous-schéma fermé défini par $g_1, \ldots, g_c$.
Notons $\pi : Y \to \mathbf{A}^{c(n + 1)}_k$ la projection
sur l'espace affine d'indéterminées $a_{ij}$.
Alors il existe un ouvert de Zariski non vide
de $\mathbf{A}^{c(n + 1)}_k$ au-dessus duquel $\pi$ est lisse.
\end{lemma}

\begin{proof}
Rappelons que l'ensemble des points où $\pi$ est lisse est ouvert.
Son complémentaire, c'est-à-dire le lieu singulier, est donc fermé.
D'après le théorème de Chevalley (sous la forme de
Morphismes, lemme \ref{morphisms-lemma-chevalley}),
l'image du lieu singulier est constructible.
Par conséquent, si le point générique de $\mathbf{A}^{c(n + 1)}_k$
n'appartient pas à l'image du lieu singulier, alors
le lemme en découle (par exemple d'après Topologie, lemme
\ref{topology-lemma-generic-point-in-constructible}).
Il faut donc montrer qu'il n'existe aucun point
$y \in Y$ où $\pi$ ne soit pas lisse et qui s'envoie sur
le point générique de $\mathbf{A}^{c(n + 1)}_k$.
Considérons la matrice des dérivées partielles
$$
(\frac{\partial g_j}{\partial x_i}) =
(\frac{\partial f_j}{\partial x_i} + a_{ij})
$$
L'image de cette matrice dans $\kappa(y)$ doit être de rang $< c$,
car sinon $\pi$ serait lisse en $y$ ; voir la discussion dans
Lissage des homomorphismes d'anneaux, section \ref{smoothing-section-singular-ideal}.
Nous pouvons donc trouver $\lambda_1, \ldots, \lambda_c \in \kappa(y)$,
non tous nuls, tels que le vecteur $(\lambda_1, \ldots, \lambda_c)$
appartienne au noyau de cette matrice.
Après renumérotation, nous pouvons supposer $\lambda_1 \not = 0$.
En divisant par $\lambda_1$, nous pouvons supposer que notre vecteur est
de la forme $(1, \lambda_2, \ldots, \lambda_c)$.
Nous obtenons alors
$$
a_{i1} = -
\frac{\partial f_j}{\partial x_1} -
\sum\nolimits_{j = 2, \ldots, c} \lambda_j(\frac{\partial f_j}{\partial x_i} + a_{ij})
$$
dans $\kappa(y)$ pour $i = 1, \ldots, n$. De plus, puisque $y \in Y$, nous avons aussi
l'égalité
$$
a_{0j} = -f_j - a_{1j}x_1 - \ldots - a_{nj}x_n
$$
dans $\kappa(y)$. Cela signifie que le sous-corps de $\kappa(y)$
engendré par $a_{ij}$ est contenu dans le sous-corps de $\kappa(y)$
engendré par les images de $x_1, \ldots, x_n, \lambda_2, \ldots, \lambda_c$,
et par les $a_{ij}$ autres que $a_{i1}$ et $a_{0j}$.
En comptant, nous voyons que son degré de transcendance est
au plus $c(n + 1) - 1$. Ainsi, $y$ ne peut pas s'envoyer sur le point générique,
comme souhaité.
\end{proof}

\begin{lemma}
\label{lemma-smoothing-affine-lci}
Soit $k$ un corps. Soit $A$ une intersection complète globale
sur $k$. Il existe un homomorphisme d'anneaux plat de type fini
$k[[t]] \to B$ tel que $B/tB \cong A$ et que
$B[1/t]$ soit lisse sur $k((t))$.
\end{lemma}

\begin{proof}
Écrivons $A = k[x_1, \ldots, x_n]/(f_1, \ldots, f_c)$ comme dans
Algèbre, définition \ref{algebra-definition-lci-field}.
Nous allons choisir
$a_{ij} \in (t) \subset k[[t]]$ et poser
$$
g_j = f_j + a_{0j} + a_{1j}x_1 + \ldots + a_{nj}x_n \in
k[[t]][x_1, \ldots, x_n]
$$
Après ce choix, nous prenons
$B = k[[t]][x_1, \ldots, x_n]/(g_1, \ldots, g_c)$.
Nous affirmons que $k[[t]] \to B$ est plat en tout idéal premier
au-dessus de $(t)$. En effet, les éléments $f_1, \ldots, f_c$
forment une suite régulière dans l'anneau local en tout idéal premier
$\mathfrak p$ de $k[x_1, \ldots, x_n]$ contenant $f_1, \ldots, f_c$
(Algèbre, lemme \ref{algebra-lemma-lci}). Ainsi, $g_1, \ldots, g_c$
est localement un relèvement d'une suite régulière, et nous pouvons appliquer
Algèbre, lemme \ref{algebra-lemma-grothendieck-regular-sequence}.
La platitude aux idéaux premiers au-dessus de $(0) \subset k[[t]]$ est automatique,
car $k((t)) = k[[t]]_{(0)}$ est un corps. Ainsi, $B$ est plat
sur $k[[t]]$.

\medskip\noindent
Il reste seulement à montrer que, pour des choix convenables
des $a_{ij}$, la fibre générique $B_{(0)}$ est lisse sur
$k((t))$. Pour cela, nous devons montrer que nous pouvons choisir
nos $a_{ij}$ de sorte que le morphisme induit
$$
(a_{ij}) : \Spec(k[[t]]) \longrightarrow \mathbf{A}^{c(n + 1)}_k
$$
s'envoie dans l'ouvert de Zariski non vide du
lemme \ref{lemma-jouanolou-type-thing}.
Cela est clair, car aucun polynôme non nul en les
$a_{ij}$ ne s'annule sur $(t)^{\oplus c(n + 1)}$.
(Nous laissons cela en exercice au lecteur.)
\end{proof}

\begin{lemma}
\label{lemma-smoothing-artinian-lci}
Soit $k$ un corps. Soit $A$ une $k$-algèbre de dimension finie
qui est une intersection complète locale sur $k$. Alors il existe
une $k[[t]]$-algèbre finie et plate $B$ telle que $B/tB \cong A$
et $B[1/t]$ soit \'etale sur $k((t))$.
\end{lemma}

\begin{proof}
Puisque $A$ est artinien
(Algèbre, lemme \ref{algebra-lemma-finite-dimensional-algebra}),
nous pouvons écrire $A$ comme un produit d'anneaux locaux artiniens
(Algèbre, lemme \ref{algebra-lemma-artinian-finite-length}).
Il suffit donc de démontrer le lemme lorsque $A$ est local
(nous utilisons ici que la propriété d'être une intersection complète locale
est préservée par localisation principale ; voir
Algèbre, lemme \ref{algebra-lemma-localize-lci}).
Dans ce cas, $A$ est une intersection complète globale.
Considérons l'algèbre $B$ construite dans le
lemme \ref{lemma-smoothing-affine-lci}.
Alors $k[[t]] \to B$ est quasi-fini en l'unique idéal premier de $B$
au-dessus de $(t)$ (Algèbre, définition \ref{algebra-definition-quasi-finite}).
Observons que $k[[t]]$ est un anneau local hensélien
(Algèbre, lemme \ref{algebra-lemma-complete-henselian}).
Ainsi, $B = B' \times C$, où $B'$ est fini sur $k[[t]]$
et $C$ n'a aucun idéal premier au-dessus de $(t)$ ; voir
Algèbre, lemme \ref{algebra-lemma-characterize-henselian}.
Alors $B'$ est l'anneau recherché
(rappelons que \'etale signifie la même chose que
lisse de dimension relative $0$).
\end{proof}

\begin{lemma}
\label{lemma-smoothing-at-lci-point}
Soit $k$ un corps. Soit $A$ une $k$-algèbre. Supposons que
\begin{enumerate}
\item $A$ soit un anneau local essentiellement de type fini sur $k$ ;
\item $A$ soit une intersection complète sur $k$
(Algèbre, définition \ref{algebra-definition-lci-local-ring}).
\end{enumerate}
Posons $d = \dim(A) + \text{trdeg}_k(\kappa)$, où $\kappa$
est le corps résiduel de $A$. Alors il existe un entier $n$
et un homomorphisme d'anneaux plat, essentiellement de type fini,
$k[[t]] \to B$ tel que $B/tB \cong A$ et que $t^n$ appartienne au
$d$-ième idéal de Fitting de $\Omega_{B/k[[t]]}$.
\end{lemma}

\begin{proof}
D'après Algèbre, lemme \ref{algebra-lemma-lci-local}, nous pouvons écrire $A$ comme la
localisation en un idéal premier $\mathfrak p$ d'une intersection complète globale $P$
sur $k$. Observons que $\dim(P) = d$ d'après
Algèbre, lemme \ref{algebra-lemma-dimension-at-a-point-finite-type-field}.
D'après le lemme \ref{lemma-smoothing-affine-lci}, nous pouvons trouver un
homomorphisme d'anneaux plat de type fini $k[[t]] \to Q$ tel que $P \cong Q/tQ$ et
tel que $k((t)) \to Q[1/t]$ soit lisse. Il résulte de la construction
de $Q$ dans le lemme que $k[[t]] \to Q$ est un morphisme
d'intersection complète globale, de dimension relative $d$ ; autrement dit,
Algèbre, lemme \ref{algebra-lemma-syntomic}, nous dit que $Q$ ou une
localisation principale convenable de $Q$ est une telle intersection complète globale.
Ainsi, d'après Diviseurs, lemme \ref{divisors-lemma-d-fitting-ideal-omega-smooth},
le $d$-ième idéal de Fitting $I \subset Q$ de $\Omega_{Q/k[[t]]}$
définit le lieu singulier de $\Spec(Q) \to \Spec(k[[t]])$.
Ainsi, $t^n \in I$ pour un certain $n$.
Soit $\mathfrak q \subset Q$
l'image réciproque de $\mathfrak p$. Posons $B = Q_\mathfrak q$.
Le lemme est démontré.
\end{proof}

\begin{lemma}
\label{lemma-smoothing-proper-curve-isolated-lci}
Soit $X$ un schéma sur un corps $k$. Supposons que
\begin{enumerate}
\item $X$ soit propre sur $k$ ;
\item $X$ soit une intersection complète locale sur $k$ ;
\item $X$ soit de dimension $\leq 1$ ;
\item $X \to \Spec(k)$ soit lisse sauf en un nombre fini de points.
\end{enumerate}
Alors il existe un morphisme projectif plat $Y \to \Spec(k[[t]])$
dont la fibre générique est lisse et la fibre spéciale
est isomorphe à $X$.
\end{lemma}

\begin{proof}
Observons que $X$ est de Cohen-Macaulay ; voir
Algèbre, lemme \ref{algebra-lemma-lci-CM}.
Ainsi, $X = X' \amalg X''$ avec $\dim(X') = 0$,
et $X''$ est équidimensionnel de dimension $1$ ; voir Morphismes, lemme
\ref{morphisms-lemma-flat-finite-presentation-CM-fibres-relative-dimension}.
Puisque $X'$ est fini sur $k$ (Variétés, lemme
\ref{varieties-lemma-algebraic-scheme-dim-0})
nous pouvons trouver $Y' \to \Spec(k[[t]])$ dont la fibre
spéciale est $X'$ et la fibre générique est lisse d'après le
lemme \ref{lemma-smoothing-artinian-lci}.
Il suffit donc de démontrer le lemme pour $X''$.
Après avoir remplacé $X$ par $X''$, le schéma $X$ est
de Cohen-Macaulay et équidimensionnel de dimension $1$.

\medskip\noindent
Nous allons utiliser la théorie des déformations dans la situation $\Lambda = k \to k$.
Soient $p_1, \ldots, p_r \in X$ les points singuliers fermés de $X$, c'est-à-dire
les points où $X \to \Spec(k)$ n'est pas lisse. Pour chaque $i$, choisissons
un entier $n_i$ et un homomorphisme d'anneaux plat, essentiellement de type fini,
$$
k[[t]] \longrightarrow B_i
$$
tel que $B_i/tB_i \cong \mathcal{O}_{X, p_i}$ et que
$t^{n_i}$ appartienne au $1$er idéal de Fitting de $\Omega_{B_i/k[[t]]}$.
Cela est possible d'après le lemme \ref{lemma-smoothing-at-lci-point}.
Observons que le système $(B_i/t^nB_i)$ définit un objet formel de
$\Deformationcategory_{\mathcal{O}_{X, p_i}}$ sur $k[[t]]$.
D'après le lemme \ref{lemma-curve-isolated}, l'application
$$
\Deformationcategory_X
\longrightarrow
\prod\nolimits_{i = 1, \ldots, r} \Deformationcategory_{\mathcal{O}_{X, p_i}}
$$
est une application lisse entre catégories de déformations. Ainsi, d'après
Théorie formelle des déformations, lemme
\ref{formal-defos-lemma-smooth-morphism-essentially-surjective}
il existe un objet formel $(X_n)$ dans $\Deformationcategory_X$
qui s'envoie sur l'objet formel $\prod_i (B_i/t^n)$ par la flèche ci-dessus.
D'après Compléments sur les morphismes d'espaces, lemme
\ref{spaces-more-morphisms-lemma-formal-algebraic-space-proper-reldim-1}
il existe un schéma projectif $Y$ sur $k[[t]]$ et des isomorphismes compatibles
$Y \times_{\Spec(k[[t]])} \Spec(k[t]/(t^n)) \cong X_n$.
D'après Compléments sur les morphismes, lemme
\ref{more-morphisms-lemma-check-flatness-on-infinitesimal-nbhds}
nous voyons que $Y \to \Spec(k[[t]])$ est plat.
Puisque $X$ est de Cohen-Macaulay et équidimensionnel de dimension $1$,
nous pouvons appliquer le lemme \ref{lemma-criterion-smoothing}
pour vérifier que $Y$ a une fibre générique lisse\footnote{Attention : en général, il n'est
{\bf pas} vrai que l'anneau local de $Y$ au point
$p_i$ soit isomorphe à $B_i$. Nous savons seulement que cela devient vrai après
quotient par $t^n$ des deux côtés !}.
Choisissons $n$ strictement supérieur au maximum des entiers $n_i$ trouvés ci-dessus.
Si nous pouvons montrer que $t^{n - 1}$ appartient au premier idéal de Fitting de
$\Omega_{X_n/S_n}$, où $S_n = \Spec(k[t]/(t^n))$, la démonstration est achevée.
Pour cela, il suffit de le démontrer dans chacun
des anneaux locaux de $X_n$ aux points fermés $p$.
Cependant, si $p$ correspond à un point lisse de $X \to \Spec(k)$,
alors $\Omega_{X_n/S_n, p}$ est libre de rang $1$ et le premier idéal de Fitting
est égal à l'anneau local. Si $p = p_i$ pour un certain $i$, alors
$$
\Omega_{X_n/S_n, p_i} =
\Omega_{(B_i/t^nB_i)/(k[t]/(t^n))} =
\Omega_{B_i/k[[t]]}/t^n\Omega_{B_i/k[[t]]}
$$
Puisque la formation des idéaux de Fitting commute au changement de base
(nous avons déjà utilisé ce fait, mais, dans ce cadre algébrique,
il résulte de Compléments sur l'algèbre, lemme
\ref{more-algebra-lemma-fitting-ideal-basics}),
et puisque $n - 1 \geq n_i$, nous voyons que $t^{n - 1}$ appartient
à l'idéal de Fitting de ce module sur $B_i/t^nB_i$, comme souhaité.
\end{proof}

\begin{lemma}
\label{lemma-smoothing-curve-isolated-lci}
Soit $k$ un corps et soit $X$ un schéma sur $k$. Supposons que
\begin{enumerate}
\item $X$ soit séparé, de type fini sur $k$ et $\dim(X) \leq 1$ ;
\item $X$ soit une intersection complète locale sur $k$ ;
\item $X \to \Spec(k)$ soit lisse sauf en un nombre fini de points.
\end{enumerate}
Alors il existe un morphisme plat, séparé et de type fini $Y \to \Spec(k[[t]])$
dont la fibre générique est lisse et la fibre spéciale
est isomorphe à $X$.
\end{lemma}

\begin{proof}
Si $X$ est réduit, nous pouvons choisir une immersion
$X \subset \overline{X}$ comme dans
Variétés, lemme \ref{varieties-lemma-reduced-dim-1-projective-completion}.
En écrivant $X = \overline{X} \setminus \{x_1, \ldots, x_n\}$,
nous voyons que $\mathcal{O}_{\overline{X}, x_i}$ est un anneau de
valuation discrète, donc en particulier une intersection complète locale
(Algèbre, définition \ref{algebra-definition-lci-local-ring}).
Ainsi, $\overline{X}$ est une intersection complète locale
sur $k$, car cette propriété vaut sur l'ouvert $X$ et
aux points $x_i$ d'après Algèbre, lemme \ref{algebra-lemma-lci-local}.
Nous pouvons donc appliquer le lemme \ref{lemma-smoothing-proper-curve-isolated-lci}
pour trouver un morphisme projectif plat $\overline{Y} \to \Spec(k[[t]])$
dont la fibre générique est lisse et la fibre spéciale
est $\overline{X}$. Nous retirons alors $x_1, \ldots, x_n$
de $\overline{Y}$ pour obtenir $Y$.

\medskip\noindent
Dans le cas général, écrivons $X = X' \amalg X''$, où
$\dim(X') = 0$ et $X''$ est équidimensionnel de dimension $1$.
Alors $X''$ est réduit et le premier paragraphe s'y applique.
D'autre part, $X'$ peut être traité
comme dans la démonstration du lemme \ref{lemma-smoothing-proper-curve-isolated-lci}.
Certains détails sont omis.
\end{proof}





\begin{multicols}{2}[\section{Autres chapitres}]
\noindent
Préliminaires
\begin{enumerate}
\item \hyperref[introduction-section-phantom]{Introduction}
\item \hyperref[conventions-section-phantom]{Conventions}
\item \hyperref[sets-section-phantom]{Théorie des ensembles}
\item \hyperref[categories-section-phantom]{Catégories}
\item \hyperref[topology-section-phantom]{Topologie}
\item \hyperref[sheaves-section-phantom]{Faisceaux sur les espaces}
\item \hyperref[sites-section-phantom]{Sites et faisceaux}
\item \hyperref[stacks-section-phantom]{Champs}
\item \hyperref[fields-section-phantom]{Corps}
\item \hyperref[algebra-section-phantom]{Algèbre commutative}
\item \hyperref[brauer-section-phantom]{Groupes de Brauer}
\item \hyperref[homology-section-phantom]{Algèbre homologique}
\item \hyperref[derived-section-phantom]{Catégories dérivées}
\item \hyperref[simplicial-section-phantom]{Méthodes simpliciales}
\item \hyperref[more-algebra-section-phantom]{Compléments d'algèbre}
\item \hyperref[smoothing-section-phantom]{Lissage des morphismes d'anneaux}
\item \hyperref[modules-section-phantom]{Faisceaux de modules}
\item \hyperref[sites-modules-section-phantom]{Modules sur les sites}
\item \hyperref[injectives-section-phantom]{Objets injectifs}
\item \hyperref[cohomology-section-phantom]{Cohomologie des faisceaux}
\item \hyperref[sites-cohomology-section-phantom]{Cohomologie sur les sites}
\item \hyperref[dga-section-phantom]{Algèbre différentielle graduée}
\item \hyperref[dpa-section-phantom]{Algèbre à puissances divisées}
\item \hyperref[sdga-section-phantom]{Faisceaux différentiels gradués}
\item \hyperref[hypercovering-section-phantom]{Hyperrecouvrements}
\end{enumerate}
Schémas
\begin{enumerate}
\setcounter{enumi}{25}
\item \hyperref[schemes-section-phantom]{Schémas}
\item \hyperref[constructions-section-phantom]{Constructions de schémas}
\item \hyperref[properties-section-phantom]{Propriétés des schémas}
\item \hyperref[morphisms-section-phantom]{Morphismes de schémas}
\item \hyperref[coherent-section-phantom]{Cohomologie des schémas}
\item \hyperref[divisors-section-phantom]{Diviseurs}
\item \hyperref[limits-section-phantom]{Limites de schémas}
\item \hyperref[varieties-section-phantom]{Variétés}
\item \hyperref[topologies-section-phantom]{Topologies sur les schémas}
\item \hyperref[descent-section-phantom]{Descente}
\item \hyperref[perfect-section-phantom]{Catégories dérivées de schémas}
\item \hyperref[more-morphisms-section-phantom]{Compléments sur les morphismes}
\item \hyperref[flat-section-phantom]{Compléments sur la platitude}
\item \hyperref[groupoids-section-phantom]{Schémas en groupoïdes}
\item \hyperref[more-groupoids-section-phantom]{Compléments sur les schémas en groupoïdes}
\item \hyperref[etale-section-phantom]{Morphismes étales de schémas}
\end{enumerate}
Thèmes de théorie des schémas
\begin{enumerate}
\setcounter{enumi}{41}
\item \hyperref[chow-section-phantom]{Homologie de Chow}
\item \hyperref[intersection-section-phantom]{Théorie de l'intersection}
\item \hyperref[pic-section-phantom]{Schémas de Picard des courbes}
\item \hyperref[weil-section-phantom]{Théories cohomologiques de Weil}
\item \hyperref[adequate-section-phantom]{Modules adéquats}
\item \hyperref[dualizing-section-phantom]{Complexes dualisants}
\item \hyperref[duality-section-phantom]{Dualité pour les schémas}
\item \hyperref[discriminant-section-phantom]{Discriminants et différentes}
\item \hyperref[derham-section-phantom]{Cohomologie de de Rham}
\item \hyperref[local-cohomology-section-phantom]{Cohomologie locale}
\item \hyperref[algebraization-section-phantom]{Géométrie algébrique et formelle}
\item \hyperref[curves-section-phantom]{Courbes algébriques}
\item \hyperref[resolve-section-phantom]{Résolution des surfaces}
\item \hyperref[models-section-phantom]{Réduction semi-stable}
\item \hyperref[functors-section-phantom]{Foncteurs et morphismes}
\item \hyperref[equiv-section-phantom]{Catégories dérivées de variétés}
\item \hyperref[pione-section-phantom]{Groupes fondamentaux de schémas}
\item \hyperref[etale-cohomology-section-phantom]{Cohomologie étale}
\item \hyperref[crystalline-section-phantom]{Cohomologie cristalline}
\item \hyperref[proetale-section-phantom]{Cohomologie pro-étale}
\item \hyperref[relative-cycles-section-phantom]{Cycles relatifs}
\item \hyperref[more-etale-section-phantom]{Compléments de cohomologie étale}
\item \hyperref[trace-section-phantom]{La formule des traces}
\end{enumerate}
Espaces algébriques
\begin{enumerate}
\setcounter{enumi}{64}
\item \hyperref[spaces-section-phantom]{Espaces algébriques}
\item \hyperref[spaces-properties-section-phantom]{Propriétés des espaces algébriques}
\item \hyperref[spaces-morphisms-section-phantom]{Morphismes d'espaces algébriques}
\item \hyperref[decent-spaces-section-phantom]{Espaces algébriques décents}
\item \hyperref[spaces-cohomology-section-phantom]{Cohomologie des espaces algébriques}
\item \hyperref[spaces-limits-section-phantom]{Limites d'espaces algébriques}
\item \hyperref[spaces-divisors-section-phantom]{Diviseurs sur les espaces algébriques}
\item \hyperref[spaces-over-fields-section-phantom]{Espaces algébriques sur des corps}
\item \hyperref[spaces-topologies-section-phantom]{Topologies sur les espaces algébriques}
\item \hyperref[spaces-descent-section-phantom]{Descente et espaces algébriques}
\item \hyperref[spaces-perfect-section-phantom]{Catégories dérivées d'espaces}
\item \hyperref[spaces-more-morphisms-section-phantom]{Compléments sur les morphismes d'espaces}
\item \hyperref[spaces-flat-section-phantom]{Platitude sur les espaces algébriques}
\item \hyperref[spaces-groupoids-section-phantom]{Groupoïdes dans les espaces algébriques}
\item \hyperref[spaces-more-groupoids-section-phantom]{Compléments sur les groupoïdes dans les espaces}
\item \hyperref[bootstrap-section-phantom]{Amorçage}
\item \hyperref[spaces-pushouts-section-phantom]{Sommes amalgamées d'espaces algébriques}
\end{enumerate}
Thèmes de géométrie
\begin{enumerate}
\setcounter{enumi}{81}
\item \hyperref[spaces-chow-section-phantom]{Groupes de Chow des espaces}
\item \hyperref[groupoids-quotients-section-phantom]{Quotients de groupoïdes}
\item \hyperref[spaces-more-cohomology-section-phantom]{Compléments sur la cohomologie des espaces}
\item \hyperref[spaces-simplicial-section-phantom]{Espaces simpliciaux}
\item \hyperref[spaces-duality-section-phantom]{Dualité pour les espaces}
\item \hyperref[formal-spaces-section-phantom]{Espaces algébriques formels}
\item \hyperref[restricted-section-phantom]{Algébrisation des espaces formels}
\item \hyperref[spaces-resolve-section-phantom]{Résolution des surfaces revisitée}
\end{enumerate}
Théorie des déformations
\begin{enumerate}
\setcounter{enumi}{89}
\item \hyperref[formal-defos-section-phantom]{Théorie formelle des déformations}
\item \hyperref[defos-section-phantom]{Théorie des déformations}
\item \hyperref[cotangent-section-phantom]{Le complexe cotangent}
\item \hyperref[examples-defos-section-phantom]{Problèmes de déformation}
\end{enumerate}
Champs algébriques
\begin{enumerate}
\setcounter{enumi}{93}
\item \hyperref[algebraic-section-phantom]{Champs algébriques}
\item \hyperref[examples-stacks-section-phantom]{Exemples de champs}
\item \hyperref[stacks-sheaves-section-phantom]{Faisceaux sur les champs algébriques}
\item \hyperref[criteria-section-phantom]{Critères de représentabilité}
\item \hyperref[artin-section-phantom]{Axiomes d'Artin}
\item \hyperref[quot-section-phantom]{Espaces de Quot et de Hilbert}
\item \hyperref[stacks-properties-section-phantom]{Propriétés des champs algébriques}
\item \hyperref[stacks-morphisms-section-phantom]{Morphismes de champs algébriques}
\item \hyperref[stacks-limits-section-phantom]{Limites de champs algébriques}
\item \hyperref[stacks-cohomology-section-phantom]{Cohomologie des champs algébriques}
\item \hyperref[stacks-perfect-section-phantom]{Catégories dérivées de champs}
\item \hyperref[stacks-introduction-section-phantom]{Introduction aux champs algébriques}
\item \hyperref[stacks-more-morphisms-section-phantom]{Compléments sur les morphismes de champs}
\item \hyperref[stacks-geometry-section-phantom]{La géométrie des champs}
\end{enumerate}
Thèmes de théorie des espaces de modules
\begin{enumerate}
\setcounter{enumi}{107}
\item \hyperref[moduli-section-phantom]{Champs de modules}
\item \hyperref[moduli-curves-section-phantom]{Espaces de modules de courbes}
\end{enumerate}
Divers
\begin{enumerate}
\setcounter{enumi}{109}
\item \hyperref[examples-section-phantom]{Exemples}
\item \hyperref[exercises-section-phantom]{Exercices}
\item \hyperref[guide-section-phantom]{Guide bibliographique}
\item \hyperref[desirables-section-phantom]{Desiderata}
\item \hyperref[coding-section-phantom]{Style de programmation}
\item \hyperref[obsolete-section-phantom]{Obsolète}
\item \hyperref[fdl-section-phantom]{Licence de documentation libre GNU}
\item \hyperref[index-section-phantom]{Index généré automatiquement}
\end{enumerate}
\end{multicols}


\bibliography{my}
\bibliographystyle{amsalpha}

\end{document}
